% chktex-file 19
% chktex-file 35
% chktex-file 36
 \documentclass[12pt]{article}
 \usepackage[left=2cm,right=2cm,top=2cm,bottom=2cm,bindingoffset=0cm]{geometry}
 \usepackage{amssymb}
 \usepackage{amsmath}
 \usepackage{amsthm}
 \usepackage{enumerate}
 \usepackage[T1,T2A]{fontenc}
 \usepackage[utf8]{inputenc}
 \usepackage[matrix,arrow,curve]{xy}
 \usepackage[colorlinks=true, urlcolor=blue, linkcolor=blue, citecolor=blue,
     pdfborder={0 0 0}]{hyperref}
\usepackage{enumitem}

 %------------------------------------------------------------------------------
 \newtheorem{theorem}{Theorem}[section]
 \newtheorem{lemma}[theorem]{Lemma}
 \newtheorem{proposition}[theorem]{Proposition}
 \newtheorem{remark}[theorem]{Remark}
 \newtheorem{corollary}[theorem]{Corollary}
 \newtheorem{definition}[theorem]{Definition}
 \newtheorem{example}[theorem]{Example}
 \newtheorem{question}[theorem]{Question}

 \newcommand{\projtens}{\mathbin{\widehat{\otimes}}}
 \newcommand{\isom}[1]{\mathop{\mathbin{\cong}}\limits_{#1}}
 \newcommand{\Nor}{\mathbf{Nor}}
 \newcommand{\Ban}{\mathbf{Ban}}
 \newcommand{\lo}[1]{\ell_1^0(#1)}
 \newcommand{\ellone}[1]{\ell_1(#1)}
 \newcommand{\norm}[1]{\lVert #1\rVert}
 \newcommand{\supp}{\operatorname{supp}}
 \newcommand{\ext}{\operatorname{ext}}
 \newcommand{\conv}{\operatorname{conv}}
 \newcommand{\Lin}{\operatorname{span}}
 \renewcommand{\Im}{\operatorname{Im}}
 \newcommand{\Ker}{\operatorname{Ker}}
 \newcommand{\BM}{d_{BM}}
 %------------------------------------------------------------------------------

\begin{document}

\begin{center}
    \Large \textbf{Extremely projective normed spaces}\\[0.5cm]
    \small {N. T. Nemesh}\\[0.5cm]
\end{center}

\thispagestyle{empty}

\medskip
\textbf{Abstract:} A normed space is called extremely projective if it is
$C$-topologically projective in the category of normed spaces for every $C>1$.
Helemskii proved that the spaces $\lo{\Lambda}$ of finitely supported functions
on a set $\Lambda$ with the $\ell_1$-norm are extremely projective and asked
whether every extremely projective normed space is isometrically isomorphic to
some $\lo{\Lambda}$; the same conjecture was raised by Gr{\o}nb{\ae}k. We answer
this question in the negative and prove that the conjecture holds in the
almost isometric form: a normed space $P$ is extremely projective if and only
if it is topologically projective and its completion is isometrically
isomorphic to $\ellone{\Gamma}$ for some set $\Gamma$, if and only if the
Banach--Mazur distance between $P$ and $\lo{\Gamma}$ equals $1$. For spaces of
countable dimension this means that $P$ is extremely projective if and only
if its completion is isometrically isomorphic to $\ell_1$ or $\ell_1^n$. The
proof is a transfinite back-and-forth construction. Along the way we show that
extremely projective normed spaces contain no infinite dimensional complete
subspaces, so that no infinite dimensional Banach space is extremely
projective in the category of normed spaces, and we use Fodor's lemma to
construct a dense subspace of $\ellone{\omega_1}$ all of whose separable
subspaces have countable dimension but which is not topologically projective.
Finally, we show that a dense subspace of $\ellone{\Gamma}$ is extremely
projective if and only if it admits a continuous filtration with countable
blocks, dense pieces and countable dimensional successive quotients. For
$|\Gamma|=\aleph_1$ this becomes a criterion in terms of a stationary set of
discontinuities, in the spirit of the theory of almost free abelian groups,
and it yields $2^{\aleph_1}$ pairwise non-isomorphic dense subspaces of
$\ellone{\omega_1}$ all of whose countably generated pieces are extremely
projective.
\medskip

\textbf{Keywords:} projective normed space, topological projectivity, extreme
projectivity, Banach--Mazur distance, $\ell_1$-space.

\bigskip

%-------------------------------------------------------------------------------
%	Introduction
%-------------------------------------------------------------------------------

\section{Introduction}\label{SectionIntroduction}

Projectivity in categories of normed and Banach modules admits several
non-equivalent versions. Besides the classical relative projectivity of
Helemskii~\cite{HelHomolBanTopAlg}, which ignores the norms of lifting
operators, there are so-called metric and topological
projectivity~\cite{HelMetrFrQMod, WhiteInjmoduAlg}, where the
norm of a lifting is controlled by a constant. In the metric theory this
constant is $1$, and in the topological theory it is some fixed $C\geq 1$. The
present paper is devoted to the limiting case between these two theories: a
module is called extremely projective if it is $C$-topologically projective for
every $C>1$. This notion was introduced and studied by
Helemskii~\cite{HelExtrProjSeqAlg, HelMetrFrQMod}, whose definition is
formally different but equivalent to ours, see
proposition~\ref{EquivWithHelemskii}. We restrict
ourselves to the simplest possible algebra, namely the zero algebra. Normed
modules over the zero algebra are just normed spaces, so we deal with extreme
projectivity in the category $\Nor$ of normed spaces and bounded linear
operators.

For Banach spaces the corresponding questions are classical. By a theorem of
K\"othe~\cite{KotheTopProjBanSp} topologically projective Banach spaces are
exactly the Banach spaces topologically isomorphic to $\ellone{\Lambda}$ for some
set $\Lambda$, and by a theorem of Grothendieck~\cite{GrothMetrProjFlatBanSp}
a Banach space has the metric lifting property if and only if it is
isometrically isomorphic to some $\ellone{\Lambda}$. For normed spaces the
answers are different, because the category $\Nor$ has its own free objects.
These are the spaces $\lo{\Lambda}$ of finitely supported functions on a set
$\Lambda$ with the $\ell_1$-norm. Gr{\o}nb{\ae}k~\cite{GronLiftProbNormSp}
extended K\"othe's theorem to $\Nor$: a normed space $E$ is topologically
projective if and only if it is topologically isomorphic to $\lo{\Lambda}$,
where $\Lambda$ is a set whose cardinality is the density character of $E$
(for infinite dimensional $E$).
Helemskii~\cite{HelMetrFrQMod} proved the metric counterpart: a normed space
is metrically projective if and only if it is isometrically isomorphic to some
$\lo{\Lambda}$. For extreme projectivity the following facts were known. The
spaces $\lo{\Lambda}$ are extremely projective, and a normed space is
extremely projective if and only if it is a near-retract of some
$\lo{\Lambda}$ \cite[proposition 2.2 and the proposition of the
introduction]{HelExtrProjSeqAlg}. Every
extremely projective normed space is, up to an isometric isomorphism, a
dense subspace of $\ellone{\Lambda}$ for some set $\Lambda$
\cite[proposition 2.3]{HelExtrProjSeqAlg}, see also
\cite[section 3]{GronLiftProbNormSp}, and every separable topologically
projective normed space has at most countable algebraic dimension
\cite[proposition 2.4]{HelExtrProjSeqAlg}. Both Helemskii
\cite[section 2]{HelExtrProjSeqAlg}, \cite[remark 3.6]{HelMetrFrQMod} and
Gr{\o}nb{\ae}k
\cite[conjecture 3.1]{GronLiftProbNormSp} asked whether, conversely, every
extremely projective normed space is isometrically isomorphic to some
$\lo{\Lambda}$; in a more detailed form, Helemskii asked which dense subspaces
of $\ellone{\Lambda}$ are extremely projective.

The main results of this paper are the following. First, the answer to the
question of Helemskii and Gr{\o}nb{\ae}k is negative: the subspace of $\ell_1$
spanned by the finitely supported sequences and a single sequence with
infinite support is extremely projective but not isometrically isomorphic to
any $\lo{\Lambda}$. Second, extreme projectivity is an almost isometric
property: a normed space $P$ is extremely projective if and only if it is
topologically projective and its completion is isometrically isomorphic to
$\ellone{\Gamma}$ for some set $\Gamma$, if and only if for every
$\varepsilon>0$ there is a linear bijection $T:P\to\lo{\Gamma}$ with
$\norm{T}\norm{T^{-1}}<1+\varepsilon$. Thus, the conjecture of Gr{\o}nb{\ae}k
holds in the almost isometric form. For spaces of countable algebraic
dimension the condition of topological projectivity is superfluous: such a
space is extremely projective if and only if its completion is isometrically
isomorphic to $\ell_1$ or to $\ell_1^n$, so every dense subspace of $\ell_1$
of countable algebraic dimension is extremely projective. Third, we extend
the dimension restriction of Helemskii to non-separable spaces: every
separable subspace of a topologically projective normed space has countable
algebraic dimension, and therefore topologically projective normed spaces
contain no infinite dimensional subspaces which are complete with respect to
a norm dominating the original one. In particular, an infinite dimensional
Banach space is never extremely projective in $\Nor$, although it may well
be extremely projective in the category $\Ban$ of Banach spaces. The last
fact explains the main technical difficulty of the subject: when the lifting
problems are posed for all normed spaces rather than Banach spaces, the
standard tools of Banach space geometry, which rely on completeness, are no
longer available, and we must work with the finite dimensional pieces of our
spaces directly. Our main tool is a back-and-forth construction, which in the
case of countable dimension replaces the decomposition method of
Pe{\l}czy\'nski in the proofs of K\"othe and Gr{\o}nb{\ae}k, and in the
general case takes the form of a transfinite recursion along a filtration of
$\Gamma$ whose perturbations are controlled on disjoint blocks of
coordinates. Fourth, in the non-separable case the dimension restriction
admits a set-theoretic refinement: by Fodor's lemma, the closed linear spans
of the initial segments of an $\omega_1$-sequence of vectors in a
topologically projective normed space must behave continuously at almost all
limit stages. Using this we construct a dense subspace of $\ellone{\omega_1}$,
built from a ladder system on $\omega_1$, which is not topologically
projective, although all its separable subspaces have countable dimension and
all its countably generated pieces are extremely projective. Thus, for
uncountable $\Lambda$ the question of Helemskii on dense subspaces of
$\ellone{\Lambda}$ does not reduce to the countable case: a dense subspace of
$\ellone{\Lambda}$ is extremely projective if and only if it is topologically
projective, and for uncountable $\Lambda$ this is a genuine restriction.
Fifth, we describe the topologically projective dense subspaces of
$\ellone{\Gamma}$ in terms of the coordinate structure of $\ellone{\Gamma}$
only: a dense subspace $P$ is extremely projective if and only if $\Gamma$
admits a continuous filtration with countable blocks along which the pieces
$P\cap\ellone{\Gamma_\delta}$ are dense, form a continuous chain and have
countable dimensional successive quotients. For $|\Gamma|=\aleph_1$ this is a
criterion in the spirit of the theory of almost free abelian groups: a dense
subspace $P$ of $\ellone{\omega_1}$ is extremely projective if and only if
its subspaces $P\cap\ellone{\delta}$, $\delta<\omega_1$, have countable
dimension and the set of limit ordinals at which the chain of these
subspaces is discontinuous is not stationary. The class of this set modulo
non-stationary sets is an isomorphism invariant of $P$, and ladder spaces
built on stationary sets provide $2^{\aleph_1}$ pairwise non-isomorphic dense
subspaces of $\ellone{\omega_1}$ all of whose countably generated pieces are
extremely projective.

Let us fix the notation and recall the definitions. All normed spaces are
considered over the complex field; the real case is identical. For a normed
space $E$ we denote by $B_E$ its closed unit ball and by $B_E^\circ$ its open
unit ball. By $\Nor$ we denote the category of normed spaces and bounded linear
operators and by $\Ban$ its full subcategory of Banach spaces. For a set
$\Lambda$ we denote by $\ellone{\Lambda}$ the Banach space of absolutely
summable families and by $\lo{\Lambda}$ its dense subspace consisting of finitely
supported families. For $\lambda\in\Lambda$ the symbol $e_\lambda$ denotes the
family which equals $1$ at $\lambda$ and $0$ elsewhere; the family
${(e_\lambda)}_{\lambda\in\Lambda}$ is a Hamel basis of $\lo{\Lambda}$. For a
normed space $E$ and a set $\Lambda$ there is a bijection between bounded
families ${(x_\lambda)}_{\lambda\in\Lambda}$ in $E$ and bounded linear
operators $T:\lo{\Lambda}\to E$, given by $T(e_\lambda)=x_\lambda$, and
$\norm{T}=\sup_{\lambda\in\Lambda}\norm{x_\lambda}$
\cite[proposition 2.1]{HelExtrProjSeqAlg}, \cite[proposition 1.10]{GronLiftProbNormSp}.
This is the sense in which the spaces $\lo{\Lambda}$ are free objects of
$\Nor$; note that no completeness of $E$ is required here, in contrast with
the case of $\ellone{\Lambda}$.

We say that a bounded linear operator
$\xi:X\to Y$ between normed spaces is $c$-topologically surjective, where
$c\geq 1$, if $c\,\xi(B_X^\circ)\supset B_Y^\circ$, and that $\xi$ is
coisometric if it is contractive and $1$-topologically surjective, that is, if
$\xi(B_X^\circ)=B_Y^\circ$. Equivalently, $\xi$ is $c$-topologically surjective
if for every $y\in Y$ and every $\delta>0$ there exists $x\in X$ with
$\xi(x)=y$ and $\norm{x}<c\norm{y}+\delta$. A linear bijection $T:E\to F$ is
called a $c$-isomorphism if $\norm{T}\norm{T^{-1}}\leq c$. The Banach--Mazur
distance between $E$ and $F$ is
\[
\BM(E,F)=\inf\{\norm{T}\norm{T^{-1}}:T:E\to F \mbox{ is a linear bijection}\},
\]
with $\BM(E,F)=+\infty$ if there are no bounded linear bijections with bounded
inverses between $E$ and $F$. A bounded linear operator $\zeta:E\to F$ is
called a $c$-retraction if there exists a bounded linear operator $\eta:F\to E$
with $\zeta\eta=1_F$ and $\norm{\zeta}\norm{\eta}\leq c$. In this case $F$ is
called a $c$-retract of $E$.

\begin{definition}\label{DefTopProj} Let $C\geq 1$. A normed space $P$ is
    called $C$-topologically projective if for every $c\geq 1$, every
    $c$-topologically surjective operator $\xi:X\to Y$ between normed spaces
    and every bounded linear operator $\phi:P\to Y$ there exists a bounded
    linear operator $\psi:P\to X$ such that $\xi\psi=\phi$ and
    $\norm{\psi}\leq cC\norm{\phi}$. A normed space is called topologically
    projective if it is $C$-topologically projective for some $C\geq 1$, and
    extremely projective if it is $C$-topologically projective for every
    $C>1$.
\end{definition}

The operator $\psi$ in definition~\ref{DefTopProj} is called a lifting of
$\phi$ along $\xi$. Extreme projectivity in the sense of
definition~\ref{DefTopProj} is precisely the extreme projectivity of
Helemskii, see proposition~\ref{EquivWithHelemskii}. We emphasize that
the spaces $X$ and $Y$ in definition~\ref{DefTopProj} are arbitrary normed
spaces. If we restrict them to Banach spaces, we obtain the notions of
topological and extreme projectivity in $\Ban$, which we shall mention only
for comparison. The two notions differ substantially, see
section~\ref{SectionNecessaryConditions}.

The paper is organized as follows. In section~\ref{SectionBasicTheory} we
collect the basic properties of extremely projective normed spaces: the
spaces $\lo{\Lambda}$ are extremely projective, extreme projectivity passes
to $(1+\varepsilon)$-retracts, and a normed space is extremely projective if
and only if it is a $(1+\varepsilon)$-retract of a free normed space for every
$\varepsilon>0$. In section~\ref{SectionNecessaryConditions} we discuss the
dimension restrictions on topologically projective normed spaces.
Section~\ref{SectionMainTheorem} is devoted to the almost isometric
characterization of extremely projective normed spaces of countable
dimension, and section~\ref{SectionCounterexample} contains the
counterexample to the question of Helemskii and Gr{\o}nb{\ae}k.
Section~\ref{SectionUncountable} treats spaces of uncountable dimension; it
contains the ladder space example and the main theorem.
Section~\ref{SectionIntrinsic} describes topologically projective dense
subspaces of $\ellone{\Gamma}$ by filtrations with countable blocks and
treats the case $|\Gamma|=\aleph_1$.

%-------------------------------------------------------------------------------
%	Basic theory
%-------------------------------------------------------------------------------

\section{Basic properties of extremely projective normed
  spaces}\label{SectionBasicTheory}

We start by comparing definition~\ref{DefTopProj} with Helemskii's one.
Our definition of extreme projectivity is equivalent to Helemskii's
definition~\cite[definition I]{HelExtrProjSeqAlg}: a normed space $P$ is
extremely projective in the sense of Helemskii if for every coisometric operator
$\xi:X\to Y$, every bounded linear operator $\phi:P\to Y$ and every
$\delta>0$ there exists a lifting $\psi:P\to X$ of $\phi$ along $\xi$ with
$\norm{\psi}<\norm{\phi}+\delta$.

\begin{proposition}\label{EquivWithHelemskii} A normed space is extremely
    projective in the sense of definition~\ref{DefTopProj} if and only if it
    is extremely projective in the sense of Helemskii.
\end{proposition}
\begin{proof} Let $P$ be extremely projective in the sense of
    definition~\ref{DefTopProj}, $\xi:X\to Y$ be a coisometric
    operator, $\phi:P\to Y$ be a bounded linear operator and $\delta>0$. If
    $\phi=0$ we take $\psi=0$. Otherwise we put
    $C=1+\delta\norm{\phi}^{-1}/2$. Since $\xi$ is $1$-topologically surjective,
    $C$-topological projectivity of $P$ gives a lifting $\psi$ of $\phi$ along
    $\xi$ with $\norm{\psi}\leq C\norm{\phi}<\norm{\phi}+\delta$.

    Conversely, let $P$ be extremely projective in the sense of Helemskii,
    $C>1$, $\xi:X\to Y$ be a
    $c$-topologically surjective operator and $\phi:P\to Y$ be a non-zero
    bounded linear operator. Consider the operator $\xi'=c\xi$ and endow the
    vector space $X$ with the norm $\norm{x}'=\max(\norm{x},\norm{\xi'(x)})$,
    which is equivalent to the original one since $\norm{x}\leq\norm{x}'\leq
    \max(1,c\norm{\xi})\norm{x}$. Denote the resulting normed space by $X'$.
    The operator $\xi':X'\to Y$ is contractive by construction. If $y\in
    B_Y^\circ$, then by $c$-topological surjectivity of $\xi$ there is $x\in
    B_X^\circ$ with $\xi'(x)=y$, and $\norm{x}'=\max(\norm{x},\norm{y})<1$.
    Thus, $\xi':X'\to Y$ is coisometric. By Helemskii's definition there is a
    lifting $\psi':P\to X'$ of $\phi$ along $\xi'$ with $\norm{\psi'}<
    \norm{\phi}+(C-1)\norm{\phi}=C\norm{\phi}$. Regarded as an operator into
    $X$, the operator $\psi=c\psi'$ satisfies $\xi\psi=\xi'\psi'=\phi$ and
    $\norm{\psi}\leq c\norm{\psi'}'\leq cC\norm{\phi}$. Since $C>1$ is arbitrary,
    $P$ is extremely projective in the sense of definition~\ref{DefTopProj}.
\end{proof}

Now we turn to the free objects of $\Nor$.

\begin{proposition}\label{FreeIsExtrProj} Let $\Lambda$ be a set. Then
    $\lo{\Lambda}$ is $(1+\varepsilon)$-topologically projective for every
    $\varepsilon>0$. In particular, $\lo{\Lambda}$ is extremely projective.
\end{proposition}
\begin{proof} By \cite[proposition 2.2]{HelExtrProjSeqAlg} the space
    $\lo{\Lambda}$ is extremely projective in the sense of Helemskii, and
    the proof of proposition~\ref{EquivWithHelemskii} shows that Helemskii's extreme
    projectivity gives exactly
    $(1+\varepsilon)$-topological projectivity for every $\varepsilon>0$.
\end{proof}

The constant $1+\varepsilon$ in proposition~\ref{FreeIsExtrProj} cannot be
replaced by $1$. In fact, with the convention of
definition~\ref{DefTopProj}, $1$-topological projectivity is an empty notion
in $\Nor$.

\begin{proposition}\label{NoOneTopProj} The only $1$-topologically projective
    normed space is the zero space.
\end{proposition}
\begin{proof} Let $X$ be the space of finitely supported sequences with the
    $\sup$-norm and $f:X\to\mathbb{C}$ be the functional
    $f(x)=s^{-1}\sum_{n=1}^\infty 2^{-n}(1-2^{-n})x_n$, where
    $s=\sum_{n=1}^\infty 2^{-n}(1-2^{-n})$. Clearly, $\norm{f}\leq 1$. The
    operator $f$ is coisometric: for $z\in B^\circ_\mathbb{C}$ choose $N$ with
    $|z|<s^{-1}\sum_{n=1}^N 2^{-n}(1-2^{-n})$, then $z=f(x)$ for the sequence
    $x=t\,\operatorname{sgn}(z)\sum_{n=1}^N e_n$ with a suitable $t\in[0,1)$,
    and $\norm{x}<1$. On the other hand, $|f(x)|=1$ for $x\in B_X$ would
    force $x_n=\operatorname{sgn}(f(x))$ for all $n$, which is impossible for
    a finitely supported sequence. Thus, $f(B_X)=f(B_X^\circ)=B^\circ_\mathbb{C}$.

    Now let $P$ be a non-zero $1$-topologically projective normed space. Take
    $p\in P$ with $\norm{p}=1$ and, by the Hahn--Banach theorem, a functional
    $g\in P^*$ with $\norm{g}=g(p)=1$. Lifting $g$ along the coisometric
    operator $f$ we get an operator $\psi:P\to X$ with $f\psi=g$ and
    $\norm{\psi}\leq 1$. Then $x=\psi(p)\in B_X$ and $f(x)=g(p)=1$, which
    contradicts the previous paragraph.
\end{proof}

Proposition~\ref{NoOneTopProj} shows why in $\Nor$ it is natural to consider
all constants $C>1$ at once, instead of the single constant $C=1$ of the
metric theory.

\begin{remark}\label{RemarkStrictConvention} The proof of
    proposition~\ref{NoOneTopProj} uses only the fact that the coisometric
    functional $f$ does not attain its norm. Functionals of this kind exist on
    Banach spaces too, for example $x\mapsto\sum_{n=1}^\infty(1-2^{-n})x_n$ on
    $\ell_1$. Therefore, if $c$-topologically surjective operators are defined
    via open balls as in definition~\ref{DefTopProj}, then $1$-topological
    projectivity is an empty notion in $\Ban$ as well, and the free Banach
    spaces $\ellone{\Lambda}$ are only $(1+\varepsilon)$-topologically projective
    for every $\varepsilon>0$. The metric theory of \cite{HelMetrFrQMod} avoids
    this problem by posing lifting problems along strictly coisometric
    operators, that is, operators mapping the closed unit ball onto the closed
    unit ball; with this convention $\ellone{\Lambda}$ and $\lo{\Lambda}$ are
    $1$-topologically projective in $\Ban$ and $\Nor$ respectively, and
    $1$-topological projectivity coincides with metric projectivity in the
    sense of \cite{HelMetrFrQMod}. For the purposes of this paper
    the choice of convention is immaterial: if we call $\xi:X\to Y$ strictly
    $c$-topologically surjective whenever $c\,\xi(B_X)\supset B_Y$ and define
    strict $C$-topological projectivity by lifting along such operators with
    the same bound $\norm{\psi}\leq cC\norm{\phi}$, then every
    $C$-topologically projective space is strictly $C$-topologically projective,
    and every strictly $C$-topologically projective space is
    $C(1+\varepsilon)$-topologically projective for every $\varepsilon>0$. The
    first claim is obvious, since strictly $c$-topologically surjective
    operators are $c$-topologically surjective. For the second one, given a
    $c$-topologically surjective operator $\xi:X\to Y$ and $\varepsilon>0$,
    endow $X$ with the equivalent norm
    $\norm{x}''=\max(\norm{x},(1+\varepsilon)c\norm{\xi(x)})$; then
    $(1+\varepsilon)c\xi$ becomes a contractive operator mapping the closed
    unit ball onto the closed unit ball (for $y\in B_Y$ the $c$-topological
    surjectivity of $\xi$ gives $x$ with $(1+\varepsilon)c\xi(x)=y$ and
    $\norm{x}<(1+\varepsilon)^{-1}\norm{y}+\varepsilon(1+\varepsilon)^{-1}\leq 1$,
    and then $\norm{x}''\leq 1$), and a lifting $\psi''$ along it with
    $\norm{\psi''}''\leq C\norm{\phi}$ gives the lifting
    $\psi=(1+\varepsilon)c\psi''$ along $\xi$ with
    $\norm{\psi}\leq(1+\varepsilon)cC\norm{\phi}$. In particular, the two
    conventions lead to the same notion of extreme projectivity.
\end{remark}

Next we record the standard permanence property of projectivity. The proof
is the same as in the Banach case, see \cite[lemma 2.7]{WhiteInjmoduAlg};
for the second part see also
\cite[proposition 1.1]{HelExtrProjSeqAlg}.

\begin{proposition}\label{RetractOfProj} Let $P'$ be a $C$-topologically
    projective normed space and $P$ be a $c$-retract of $P'$. Then $P$ is
    $cC$-topologically projective. Consequently, a normed space which is a
    $(1+\varepsilon)$-retract of an extremely projective normed space for every
    $\varepsilon>0$ is extremely projective.
\end{proposition}

For a normed space $P$ consider the operator
\[
\pi_P:\lo{B_P}\to P,\quad e_x\mapsto x.
\]
It is well defined and contractive, because the family
${(x)}_{x\in B_P}$ is bounded by $1$. Moreover, $\pi_P$ maps the closed unit
ball of $\lo{B_P}$ onto $B_P$, since $\pi_P(e_x)=x$ for $x\in B_P$. Thus,
$\pi_P$ is strictly coisometric and, a fortiori, coisometric.

\begin{proposition}\label{CanonicalMorphismCriterion} Let $P$ be a normed
    space and $C\geq 1$. Then
    \begin{enumerate}[label = (\roman*)]
        \item if $P$ is $C$-topologically projective, then $\pi_P$ has a right
              inverse bounded linear operator of norm at most $C$;
        \item if $\pi_P$ has a right inverse bounded linear operator of norm at
              most $C$, then $P$ is $C(1+\varepsilon)$-topologically projective
              for every $\varepsilon>0$.
    \end{enumerate}
\end{proposition}
\begin{proof} $(i)$ Lift $1_P$ along the $1$-topologically surjective operator
    $\pi_P$: we get $\sigma:P\to\lo{B_P}$ with $\pi_P\sigma=1_P$ and
    $\norm{\sigma}\leq C$.

    $(ii)$ If $\pi_P\sigma=1_P$ with $\norm{\sigma}\leq C$, then $P$ is a
    $C$-retract of $\lo{B_P}$, which is $(1+\varepsilon)$-topologically
    projective by proposition~\ref{FreeIsExtrProj}. It remains to apply
    proposition~\ref{RetractOfProj}.
\end{proof}

The equivalence of conditions $(i)$ and $(iii)$ in the following theorem is
the main observation of \cite[section 2]{HelExtrProjSeqAlg} on extremely
projective normed spaces: a normed space is extremely projective if and only
if it is a near-retract of a free normed space, where a near-retract is a
retract via a contractive operator which has right inverses of norm
$<1+\varepsilon$ for every $\varepsilon>0$.

\begin{theorem}\label{ExtrProjCriterion} For a normed space $P$ the following
    conditions are equivalent:
    \begin{enumerate}[label = (\roman*)]
        \item $P$ is extremely projective;
        \item for every $\varepsilon>0$ the operator $\pi_P$ has a right
              inverse bounded linear operator of norm at most $1+\varepsilon$;
        \item for every $\varepsilon>0$ there is a set $\Lambda$ such that $P$
              is a $(1+\varepsilon)$-retract of $\lo{\Lambda}$;
        \item for every $\varepsilon>0$ there are a set $\Lambda$, a subspace
              $Z\subset\lo{\Lambda}$ and a projection $Q:\lo{\Lambda}\to\lo{\Lambda}$
              onto $Z$ with $\norm{Q}\leq 1+\varepsilon$ such that
              $\BM(P,Z)\leq 1+\varepsilon$.
    \end{enumerate}
\end{theorem}
\begin{proof} $(i)\implies(ii)$ is proposition~\ref{CanonicalMorphismCriterion}~$(i)$,
    and $(ii)\implies(iii)$ is obvious.

    $(iii)\implies(iv)$ Let $\zeta:\lo{\Lambda}\to P$ and $\eta:P\to\lo{\Lambda}$
    satisfy $\zeta\eta=1_P$ and $\norm{\zeta}\norm{\eta}\leq 1+\varepsilon$. Replacing
    $\zeta$ by $\norm{\zeta}^{-1}\zeta$ and $\eta$ by $\norm{\zeta}\eta$ we may
    assume that $\norm{\zeta}=1$ and $\norm{\eta}\leq 1+\varepsilon$. Put
    $Z=\eta(P)$ and $Q=\eta\zeta$. Then $Q^2=\eta(\zeta\eta)\zeta=Q$,
    $\Im(Q)=Z$ and $\norm{Q}\leq 1+\varepsilon$. The operator $\eta:P\to Z$
    is a bijection with $\norm{\eta}\leq 1+\varepsilon$ and
    $\eta^{-1}=\zeta|_Z$, so $\norm{\eta^{-1}}\leq 1$ and
    $\BM(P,Z)\leq 1+\varepsilon$.

    $(iv)\implies(i)$ Since $\BM(P,Z)\leq 1+\varepsilon$, there is a linear
    bijection $T:P\to Z$ with $\norm{T}\norm{T^{-1}}\leq 1+2\varepsilon$. Put
    $\eta=T$ and $\zeta=T^{-1}Q:\lo{\Lambda}\to P$. Then
    $\zeta\eta=T^{-1}QT=T^{-1}T=1_P$ and
    $\norm{\zeta}\norm{\eta}\leq(1+2\varepsilon)(1+\varepsilon)$. By
    propositions~\ref{FreeIsExtrProj} and~\ref{RetractOfProj} the space $P$ is
    $(1+2\varepsilon)(1+\varepsilon)^2$-topologically projective. As
    $\varepsilon>0$ is arbitrary, $P$ is extremely projective.
\end{proof}

%-------------------------------------------------------------------------------
%	Necessary conditions
%-------------------------------------------------------------------------------

\section{Dimension restrictions on topologically projective normed
  spaces}\label{SectionNecessaryConditions}

The free normed spaces $\lo{\Lambda}$ are far from being complete: for an
infinite set $\Lambda$ the Hamel basis ${(e_\lambda)}_{\lambda\in\Lambda}$ of
$\lo{\Lambda}$ has the same cardinality as a dense subset of $\lo{\Lambda}$
of minimal cardinality, whereas the algebraic dimension of an infinite
dimensional Banach space is uncountable by the Baire category theorem, and in
fact at least continuum~\cite{LowigDimLinRaume}. In this section we show that
this smallness is inherited by all topologically projective normed spaces.
Helemskii proved that a separable topologically projective normed space has at
most countable algebraic dimension \cite[proposition 2.4]{HelExtrProjSeqAlg},
and deduced that $\ellone{\Lambda}$ is not topologically projective in $\Nor$
for any infinite set $\Lambda$ \cite[proposition 2.5]{HelExtrProjSeqAlg}. His
argument applies verbatim to subspaces of arbitrary density character, and in
this hereditary form it will be convenient for us.

\begin{theorem}\label{SeparableSubspaces} Let $P$ be a topologically
    projective normed space and $Z$ be a subspace of $P$. Then the algebraic
    dimension of $Z$ is at most $\max(\operatorname{dens}(Z),\aleph_0)$,
    where $\operatorname{dens}(Z)$ is the density character of $Z$, that is,
    the smallest cardinality of a dense subset of $Z$. In particular, every
    separable subspace of $P$ has at most countable algebraic dimension.
\end{theorem}
\begin{proof} By proposition~\ref{CanonicalMorphismCriterion}~$(i)$ there
    exists a bounded linear operator $\sigma:P\to\lo{B_P}$ with
    $\pi_P\sigma=1_P$; in particular, $\sigma$ is injective. Let $M$ be a
    dense subset of $Z$ of cardinality $\operatorname{dens}(Z)$ and
    $S_0=\bigcup_{x\in M}\supp(\sigma(x))$; since the supports are finite,
    $S_0$ has cardinality at most $\max(\operatorname{dens}(Z),\aleph_0)$. If
    $z\in Z$, then $\sigma(z)$ is the limit in $\ellone{B_P}$ of a sequence of
    elements of $\sigma(M)$, which are supported in $S_0$; since the coordinate
    functionals are continuous, $\sigma(z)$ is supported in $S_0$ as well.
    Thus, $\sigma(Z)\subset\lo{S_0}=\Lin\{e_s:s\in S_0\}$, a space whose
    algebraic dimension is the cardinality of $S_0$, and $\sigma|_Z$ is
    injective.
\end{proof}

\begin{corollary}\label{NoCompleteSubspaces} Let $P$ be a topologically
    projective normed space, $Z$ be a subspace of $P$ and $\norm{\cdot}'$ be a
    norm on $Z$ such that $(Z,\norm{\cdot}')$ is complete and
    $\norm{z}\leq K\norm{z}'$ for all $z\in Z$ and some constant $K>0$. Then
    $Z$ is finite dimensional. In particular, every complete subspace of $P$
    is finite dimensional, and no infinite dimensional Banach space is
    topologically projective in $\Nor$ (for the spaces $\ellone{\Lambda}$ this
    is \cite[proposition 2.5]{HelExtrProjSeqAlg}).
\end{corollary}
\begin{proof} Assume that $Z$ is infinite dimensional. Choose a linearly
    independent sequence ${(z_n)}_{n\in\mathbb{N}}$ in $Z$ and let $Z_0$ be
    the closure of its linear span in $(Z,\norm{\cdot}')$. Then
    $(Z_0,\norm{\cdot}')$ is a separable infinite dimensional Banach space, so
    its algebraic dimension is uncountable. On the other hand, $Z_0$ is
    separable with respect to $\norm{\cdot}$ because $\norm{\cdot}\leq
    K\norm{\cdot}'$, so by theorem~\ref{SeparableSubspaces} the algebraic
    dimension of $Z_0$ is countable. This contradiction proves the first
    claim; the rest follows by taking $Z=P$ and $\norm{\cdot}'=\norm{\cdot}$.
\end{proof}

Let us compare theorem~\ref{SeparableSubspaces} with the theorem of
Gr{\o}nb{\ae}k \cite[theorem 2.1]{GronLiftProbNormSp}: a normed space $E$ is
topologically projective if and only if it is topologically isomorphic to
$\lo{\Lambda}$, where the cardinality of $\Lambda$ is the density character
of $E$. Since for infinite $\Lambda$ the algebraic dimension of $\lo{\Lambda}$
equals the cardinality of $\Lambda$, the algebraic dimension of an infinite
dimensional topologically projective normed space coincides with its density
character. Theorem~\ref{SeparableSubspaces} is a hereditary version of this
equality: it applies to all subspaces of $P$, which are, of course, not
topologically projective in general.

The following example shows that corollary~\ref{NoCompleteSubspaces} is
sharper than its last claim: a normed space may have no infinite dimensional
complete subspaces and still fail to be topologically projective. It also
shows that topological projectivity of a normed space is not determined by its
completion.

\begin{example}\label{WeightedExample} Let ${(w_n)}_{n\in\mathbb{N}}$ be a
    sequence of real numbers with $w_n\geq 1$ and $\lim_n w_n=+\infty$, and let
    \[
    P_w=\Bigl\{x\in\ell_1:\sum_{n=1}^\infty w_n|x_n|<+\infty\Bigr\}
    \]
    with the norm inherited from $\ell_1$. Then $P_w$ is a dense subspace of
    $\ell_1$ containing $\lo{\mathbb{N}}$, so its completion is $\ell_1$, and
    \begin{enumerate}[label = (\roman*)]
        \item every complete subspace of $P_w$ is finite dimensional;
        \item $P_w$ is not topologically projective.
    \end{enumerate}
\end{example}
\begin{proof} Denote $\norm{x}_w=\sum_{n=1}^\infty w_n|x_n|$ for $x\in P_w$.
    Then $(P_w,\norm{\cdot}_w)$ is the Banach space $\ell_1(w)$ of
    $w$-weighted summable sequences and $\norm{x}\leq\norm{x}_w$ for $x\in
    P_w$.

    $(i)$ Let $Z$ be a complete subspace of $P_w$. If ${(z_k)}_{k\in\mathbb{N}}$
    is a sequence in $Z$ converging to $z\in P_w$ with respect to
    $\norm{\cdot}_w$, then it converges to $z$ in $\ell_1$, so $z\in Z$. Thus,
    $Z$ is closed in $(P_w,\norm{\cdot}_w)$ and the two norms $\norm{\cdot}$
    and $\norm{\cdot}_w$ are complete on $Z$. Since $\norm{\cdot}\leq
    \norm{\cdot}_w$, by the open mapping theorem there is a constant $K$ with
    $\norm{z}_w\leq K\norm{z}$ for all $z\in Z$. Assume that $Z$ is infinite
    dimensional and fix $N\in\mathbb{N}$. The linear map $Z\to\mathbb{C}^N$,
    $z\mapsto(z_1,\ldots,z_N)$ has a non-zero kernel, so there exists $z\in
    Z$ with $\norm{z}=1$ and $z_n=0$ for $n\leq N$. Then
    $K\geq\norm{z}_w=\sum_{n>N}w_n|z_n|\geq\inf_{n>N}w_n$. Since $N$ is
    arbitrary and $\lim_n w_n=+\infty$, we get a contradiction.

    $(ii)$ Apply corollary~\ref{NoCompleteSubspaces} to $Z=P_w$ and
    $\norm{\cdot}'=\norm{\cdot}_w$. Alternatively, note that $P_w$ is
    separable, being dense in $\ell_1$, while its algebraic dimension is
    uncountable, since $(P_w,\norm{\cdot}_w)$ is an infinite dimensional
    Banach space; so theorem~\ref{SeparableSubspaces} applies directly.
\end{proof}

%-------------------------------------------------------------------------------
%	Main theorem
%-------------------------------------------------------------------------------

\section{Extremely projective normed spaces of countable
  dimension}\label{SectionMainTheorem}

We start with a necessary condition which holds in all dimensions. Recall that
the completion $\overline{P}$ of a normed space $P$ is a Banach space containing
$P$ as a dense subspace. Every bounded linear operator between normed spaces
extends uniquely to a bounded linear operator between their completions with
the same norm, and these extensions respect compositions. The following
theorem is \cite[proposition 2.3]{HelExtrProjSeqAlg}; see also
\cite[section 3]{GronLiftProbNormSp}.

\begin{theorem}\label{CompletionIsEllOne} Let $P$ be an extremely projective
    normed space. Then $\overline{P}$ is isometrically isomorphic to
    $\ellone{\Gamma}$ for some set $\Gamma$.
\end{theorem}
\begin{proof} We only indicate the argument. By theorem~\ref{ExtrProjCriterion}
    for every $\varepsilon>0$ the space $\overline{P}$ is a
    $(1+\varepsilon)$-retract of the Banach space $\ellone{B_P}$, so by
    Grothendieck's theorem \cite{GrothMetrProjFlatBanSp} in the form given
    in~\cite[section 2]{HelExtrProjSeqAlg} (see also \cite[proposition 3.2]{HelMetrFrQMod})
    $\overline{P}$ is isometrically isomorphic to $\ellone{\Gamma}$ for some
    set $\Gamma$. Explicitly, let $\overline{\pi}:\ellone{B_P}\to\overline{P}$
    be the extension of $\pi_P$ and $\overline{\sigma}_\varepsilon$ be the
    extensions of right inverses of $\pi_P$ of norm at most $1+\varepsilon$.
    Taking adjoints one sees that ${\overline{\pi}}^*$ is an isometric
    weak$^*$-continuous embedding of ${\overline{P}}^*$ into
    $\ell_\infty(B_P)$ with weak$^*$-closed range and that
    ${\overline{\pi}}^*{\overline{\sigma}_\varepsilon}^*$ are projections onto
    this range of norm at most $1+\varepsilon$; a limit point of these
    projections in the weak$^*$ operator topology is a projection of norm
    $1$, so ${\overline{P}}^*$ is a $1$-complemented subspace of the
    $1$-injective space $\ell_\infty(B_P)$. Hence ${\overline{P}}^*$ is
    $1$-injective, hence isometrically isomorphic to $C(K)$ for a Stonean
    space $K$ \cite[chapter 3, \S 11, theorem 6]{LaceyIsomThOfClassicBanSp}; a Banach space
    whose dual is $1$-injective is isometrically isomorphic to some
    $L_1(\Omega,\mu)$ \cite[chapter 7]{LaceyIsomThOfClassicBanSp}; and
    $\overline{P}$ has the Schur property as a subspace of
    $\ellone{B_P}$, which forces $\mu$ to be purely atomic.
\end{proof}

Theorems~\ref{SeparableSubspaces}
and~\ref{CompletionIsEllOne} describe extremely projective normed spaces
from two opposite sides: they are dense subspaces of the spaces
$\ellone{\Gamma}$ which are small enough to have separable subspaces of
countable algebraic dimension only. In this section we show that for spaces of
countable dimension there is nothing in between: every dense subspace of
$\ell_1$ of countable algebraic dimension is extremely projective, and even
almost isometric to $\lo{\mathbb{N}}$. This answers the question of
Helemskii~\cite[section 2]{HelExtrProjSeqAlg} on dense subspaces of
$\ellone{\Lambda}$ for countable $\Lambda$. The proof is a back-and-forth
construction based on the following three lemmas. For a sequence
$x\in\ell_1$ and $N\in\mathbb{N}\cup\{0\}$ we denote by $x|_N$ the sequence
which coincides with $x$ at $n\leq N$ and vanishes at $n>N$.

\begin{lemma}\label{OneVectorTheorem} Let $v\in\ell_1$,
    $N\in\mathbb{N}\cup\{0\}$ and $0<\delta<1$. Then there exists a bounded
    linear operator $R:\ell_1\to\ell_1$ with $\norm{R}\leq\delta$ such that
    \begin{enumerate}[label = (\roman*)]
        \item $R(x)=0$ for every $x\in\ell_1$ with $x=x|_N$;
        \item $R(\lo{\mathbb{N}})\subset\lo{\mathbb{N}}$;
        \item $(1_{\ell_1}-R)(v)\in\lo{\mathbb{N}}$;
        \item ${(1_{\ell_1}-R)}^{-1}(\lo{\mathbb{N}})\subset\lo{\mathbb{N}}+\mathbb{C}v$.
    \end{enumerate}
    In particular, $T=1_{\ell_1}-R$ is a linear bijection of $\ell_1$ with
    $\norm{T}\leq 1+\delta$ and $\norm{T^{-1}}\leq{(1-\delta)}^{-1}$ such
    that $T(x)=x$ whenever $x=x|_N$, $T(\lo{\mathbb{N}})\subset\lo{\mathbb{N}}$
    and $T(\lo{\mathbb{N}}+\mathbb{C}v)=\lo{\mathbb{N}}$.
\end{lemma}
\begin{proof} Since $\norm{R}\leq\delta<1$, the operator $T=1_{\ell_1}-R$ is a
    linear bijection with $T^{-1}=\sum_{m=0}^\infty R^m$,
    $\norm{T}\leq 1+\delta$ and $\norm{T^{-1}}\leq{(1-\delta)}^{-1}$, and the
    properties of $T$ listed in the last sentence follow from $(i)$--$(iv)$.
    So we only need to construct $R$ with the properties $(i)$--$(iv)$.

    If $v-v|_N$ is finitely supported, take $R=0$. Assume that $v-v|_N$ has
    infinite support. For $M\geq N$ put $t_M=\norm{v-v|_M}=\sum_{n>M}|v_n|$.
    Then the sequence ${(t_M)}_{M\geq N}$ is non-increasing, $t_M>0$ for all
    $M\geq N$ because $v-v|_N$ has infinite support, and $\lim_M t_M=0$.
    Choose integers $N=N_0<N_1<N_2<\ldots$ such that
    $t_{N_k}\leq(\delta/2)t_{N_{k-1}}$ for all $k\in\mathbb{N}$. The sets
    $B_k=\{n\in\mathbb{N}:N_{k-1}<n\leq N_k\}$, $k\in\mathbb{N}$, are finite,
    pairwise disjoint and cover $\{n\in\mathbb{N}:n>N\}$. Put
    $w_k=v|_{N_k}-v|_{N_{k-1}}$, that is, $w_k$ is the restriction of $v$ to
    $B_k$. Then $w_k\in\lo{\mathbb{N}}$, $v-v|_N=\sum_{k=1}^\infty w_k$ in
    $\ell_1$, $\supp(w_k)\subset B_k$, and
    \[
    \norm{w_k}=t_{N_{k-1}}-t_{N_k}\geq(1-\delta/2)t_{N_{k-1}}>0,\qquad
    \norm{w_{k+1}}\leq t_{N_k}\leq(\delta/2)t_{N_{k-1}}\leq
    \frac{\delta/2}{1-\delta/2}\norm{w_k}\leq\delta\norm{w_k},
    \]
    where the last inequality holds because $\delta\leq 1$. For
    $k\in\mathbb{N}$ define the functional
    \[
    g_k:\ell_1\to\mathbb{C},\quad
    g_k(x)=\sum_{n\in B_k}\overline{\operatorname{sgn}(v_n)}\,x_n,
    \]
    where $\operatorname{sgn}(z)=z/|z|$ for $z\neq 0$ and
    $\operatorname{sgn}(0)=0$. Then $|g_k(x)|\leq\sum_{n\in B_k}|x_n|$ for
    all $x\in\ell_1$, $g_k(w_k)=\norm{w_k}$ and $g_k(w_j)=0$ for $j\neq k$,
    because the supports of $w_j$ and $g_k$ are disjoint. Now define
    \[
    R:\ell_1\to\ell_1,\quad
    R(x)=\sum_{k=1}^\infty\frac{g_k(x)}{\norm{w_k}}\,w_{k+1}.
    \]
    The series converges absolutely, since
    \[
    \sum_{k=1}^\infty\frac{|g_k(x)|}{\norm{w_k}}\norm{w_{k+1}}
    \leq\delta\sum_{k=1}^\infty\sum_{n\in B_k}|x_n|\leq\delta\norm{x},
    \]
    so $R$ is a well defined bounded linear operator with $\norm{R}\leq\delta$.
    We verify the properties $(i)$--$(iv)$.

    $(i)$ If $x=x|_N$, then $g_k(x)=0$ for all $k$, because
    $\supp(g_k)\subset B_k\subset\{n:n>N\}$. Hence $R(x)=0$.

    $(ii)$ If $x\in\lo{\mathbb{N}}$, then $g_k(x)\neq 0$ only for the finitely
    many $k$ with $B_k\cap\supp(x)\neq\varnothing$, so $R(x)$ is a finite
    linear combination of the finitely supported sequences $w_{k+1}$.

    $(iii)$ Since $g_k(v|_N)=0$ and
    $g_k(v-v|_N)=\sum_{j=1}^\infty g_k(w_j)=g_k(w_k)=\norm{w_k}$, we get
    $R(v)=\sum_{k=1}^\infty w_{k+1}=v-v|_N-w_1$, so
    $(1_{\ell_1}-R)(v)=v|_N+w_1\in\lo{\mathbb{N}}$.

    $(iv)$ Let $x\in\lo{\mathbb{N}}$ and let $K$ be the finite set of those
    $k$ for which $g_k(x)\neq 0$. Put $c_k=g_k(x)/\norm{w_k}$ for $k\in K$.
    Since $g_j(w_k)=\delta_{jk}\norm{w_k}$, we have $R(w_k)=w_{k+1}$ for all
    $k$, hence $R^m(x)=\sum_{k\in K}c_k w_{k+m}$ for $m\geq 1$ and
    \[
    {(1_{\ell_1}-R)}^{-1}(x)=\sum_{m=0}^\infty R^m(x)
    =x+\sum_{k\in K}c_k\sum_{m=1}^\infty w_{k+m}
    =x+\sum_{k\in K}c_k(v-v|_{N_k})\in\lo{\mathbb{N}}+\mathbb{C}v.
    \]
\end{proof}

\begin{lemma}\label{HittingLemma} Let $Y$ be a finite dimensional subspace
    of $\lo{\mathbb{N}}$ and $n\in\mathbb{N}$ be such that $e_n\notin Y$.
    Then there is a constant $C\geq 1$ such that for every
    $z\in\lo{\mathbb{N}}$ with $\norm{z-e_n}\leq{(4C)}^{-1}$ there exists a
    linear bijection $S:\ell_1\to\ell_1$ with
    $\norm{S-1_{\ell_1}}\leq 2C\norm{z-e_n}\leq 1/2$,
    $S(y)=y$ for all $y\in Y$ and $S(z)=e_n$.
\end{lemma}
\begin{proof} Since $Y$ is finite dimensional, it is closed in $\ell_1$, so
    $d=\inf\{\norm{e_n-y}:y\in Y\}>0$. By the Hahn--Banach theorem there is
    $f\in\ell_1^*$ with $f(y)=0$ for $y\in Y$, $f(e_n)=1$ and $\norm{f}=d^{-1}$.
    Put $C=\max(1,d^{-1})$. Let $z\in\lo{\mathbb{N}}$ satisfy
    $\norm{z-e_n}\leq{(4C)}^{-1}$. Then
    $|f(z)-1|=|f(z-e_n)|\leq C\norm{z-e_n}\leq 1/4$, so $|f(z)|\geq 3/4$.
    Define
    \[
    S:\ell_1\to\ell_1,\quad S(x)=x+\frac{f(x)}{f(z)}(e_n-z).
    \]
    Then $S(y)=y$ for $y\in Y$ because $f$ vanishes on $Y$, and
    $S(z)=z+(e_n-z)=e_n$. Moreover,
    $\norm{S-1_{\ell_1}}\leq\norm{f}|f(z)|^{-1}\norm{e_n-z}\leq
    \frac43 C\norm{z-e_n}\leq 2C\norm{z-e_n}\leq 1/2$, so $S$ is a linear
    bijection of $\ell_1$ by the Neumann series.
\end{proof}

\begin{lemma}\label{ProductLemma} Let $X$ be a Banach space and
    ${(F_m)}_{m\in\mathbb{N}}$ be a sequence of bounded linear operators on $X$
    such that $\theta_m=\norm{F_m-1_X}<1$ for all $m$ and
    $\sum_{m=1}^\infty\theta_m<+\infty$. Put $A_j=F_jF_{j-1}\cdots F_1$ for
    $j\in\mathbb{N}$. Then the sequence ${(A_j)}_{j\in\mathbb{N}}$ converges in
    the operator norm to a linear bijection $T$ of $X$ with
    \[
    \norm{T}\leq\prod_{m=1}^\infty(1+\theta_m),\qquad
    \norm{T^{-1}}\leq\prod_{m=1}^\infty{(1-\theta_m)}^{-1}.
    \]
\end{lemma}
\begin{proof} Both products converge because $\sum_m\theta_m<+\infty$; denote
    them by $M$ and $M'$. Then $\norm{A_j}\leq M$ and, since each $F_m$ is
    invertible with $\norm{F_m^{-1}}\leq{(1-\theta_m)}^{-1}$ by the Neumann
    series, $\norm{A_j^{-1}}\leq M'$ for all $j$. From
    $\norm{A_j-A_{j-1}}=\norm{(F_j-1_X)A_{j-1}}\leq\theta_j M$ we see that
    ${(A_j)}_j$ is a Cauchy sequence; let $T$ be its limit, so $\norm{T}\leq M$.
    Also $\norm{A_j^{-1}-A_k^{-1}}=\norm{A_j^{-1}(A_k-A_j)A_k^{-1}}\leq
    {M'}^2\norm{A_k-A_j}$, so ${(A_j^{-1})}_j$ is a Cauchy sequence with a
    limit $B$, $\norm{B}\leq M'$. Passing to the limit in
    $A_jA_j^{-1}=A_j^{-1}A_j=1_X$ we get $TB=BT=1_X$.
\end{proof}

\begin{theorem}\label{BackAndForth} Let $D$ be a dense subspace of $\ell_1$ of
    countable algebraic dimension and $\varepsilon>0$. Then there exists a
    linear bijection $T:\ell_1\to\ell_1$ with
    $\norm{T}\norm{T^{-1}}\leq 1+\varepsilon$ and $T(D)=\lo{\mathbb{N}}$. In
    particular, $\BM(D,\lo{\mathbb{N}})=1$.
\end{theorem}
\begin{proof} Since $\ell_1$ is infinite dimensional and $D$ is dense in
    $\ell_1$, the algebraic dimension of $D$ is countably infinite. Let
    ${(d_j)}_{j\in\mathbb{N}}$ be a Hamel basis of $D$. Choose a sequence
    ${(\theta_m)}_{m\in\mathbb{N}}$ of numbers in $(0,1/2]$ with
    $\prod_{m=1}^\infty(1+\theta_m){(1-\theta_m)}^{-1}\leq 1+\varepsilon$;
    then $\sum_m\theta_m<+\infty$, since $(1+\theta)/(1-\theta)\geq 1+2\theta$.

    We construct by induction linear bijections $F_1,F_2,\ldots$ of $\ell_1$
    with $\norm{F_m-1_{\ell_1}}\leq\theta_m$ and finite dimensional subspaces
    $Y_0=\{0\}\subset Y_1\subset Y_2\subset\ldots$ of $\lo{\mathbb{N}}$ such
    that, with $A_0=1_{\ell_1}$ and $A_j=F_{2j}F_{2j-1}A_{j-1}$, for every
    $j\in\mathbb{N}$ the following conditions hold:
    \begin{enumerate}[label = (\alph*)]
        \item $A_j(d_i)\in Y_j$ for $i\leq j$;
        \item $e_i\in Y_j$ for $i\leq j$;
        \item $Y_j\subset A_j(D)$;
        \item $F_{2j-1}$ and $F_{2j}$ act as the identity on $Y_{j-1}$.
    \end{enumerate}
    Suppose that $A_{j-1}$ and $Y_{j-1}$ have been constructed (for $j=1$
    there is nothing to suppose). Let $N$ be an integer not smaller than every
    element of the supports of the vectors of some basis of $Y_{j-1}$, so that
    $y=y|_N$ for all $y\in Y_{j-1}$. Apply lemma~\ref{OneVectorTheorem} to
    the vector $v=A_{j-1}(d_j)$, the integer $N$ and $\delta=\theta_{2j-1}$,
    and let $F_{2j-1}=1_{\ell_1}-R$ be the resulting operator. Then
    $\norm{F_{2j-1}-1_{\ell_1}}\leq\theta_{2j-1}$, $F_{2j-1}$ acts as the
    identity on $Y_{j-1}$, and $y_j=F_{2j-1}A_{j-1}(d_j)\in\lo{\mathbb{N}}$.
    Put $A'=F_{2j-1}A_{j-1}$ and $Y'=Y_{j-1}+\mathbb{C}y_j\subset\lo{\mathbb{N}}$.
    Note that $Y'\subset A'(D)$: indeed, $Y_{j-1}\subset A_{j-1}(D)$ by the
    induction hypothesis (for $j=1$ trivially), $F_{2j-1}$ fixes $Y_{j-1}$, so
    $Y_{j-1}=F_{2j-1}(Y_{j-1})\subset A'(D)$, and $y_j=A'(d_j)\in A'(D)$.

    If $e_j\in Y'$, put $F_{2j}=1_{\ell_1}$ and $Y_j=Y'$. Otherwise, let $C$
    be the constant given by lemma~\ref{HittingLemma} for $Y'$ and $n=j$, and
    let $N'$ be an integer not smaller than $j$ and than every element of the
    supports of the vectors of some basis of $Y'$. Since $A'$ is a linear
    bijection of $\ell_1$ and $D$ is dense, $A'(D)$ is dense in $\ell_1$, so
    there is $d\in D$ with
    \[
    \norm{A'(d)-e_j}\leq\frac{\theta_{2j}}{16C}.
    \]
    Apply lemma~\ref{OneVectorTheorem} to the vector $A'(d)$, the integer
    $N'$ and $\delta=\theta_{2j}/2$, and let $G=1_{\ell_1}-R'$ be the
    resulting operator. Then $G$ fixes $Y'$ and $e_j$,
    $\norm{G-1_{\ell_1}}\leq\theta_{2j}/2$, $\norm{G}\leq 5/4$ because
    $\theta_{2j}\leq 1/2$, and $z=GA'(d)\in\lo{\mathbb{N}}$ satisfies
    \[
    \norm{z-e_j}=\norm{G(A'(d)-e_j)}\leq\frac{5\theta_{2j}}{64C}\leq\frac{1}{4C}.
    \]
    Let $S$ be the operator given by lemma~\ref{HittingLemma} for $Y'$, $n=j$
    and $z$. Then $S$ fixes $Y'$, $S(z)=e_j$ and
    $\norm{S-1_{\ell_1}}\leq 2C\norm{z-e_j}\leq 5\theta_{2j}/32$. Put $F_{2j}=SG$
    and $Y_j=Y'+\mathbb{C}e_j$. Then
    \[
    \norm{F_{2j}-1_{\ell_1}}\leq\norm{(S-1_{\ell_1})G}+\norm{G-1_{\ell_1}}
    \leq\frac{5\theta_{2j}}{32}\cdot\frac54+\frac{\theta_{2j}}{2}
    \leq\theta_{2j}.
    \]
    Let us check $(a)$--$(d)$ for $j$. The operator $F_{2j}$ acts as the
    identity on $Y'\supset Y_{j-1}$, which together with the first step gives
    $(d)$. For $i<j$ we have $A_{j-1}(d_i)\in Y_{j-1}$, so
    $A_j(d_i)=F_{2j}F_{2j-1}A_{j-1}(d_i)=A_{j-1}(d_i)\in Y_j$, and
    $A_j(d_j)=F_{2j}(y_j)=y_j\in Y_j$; this is $(a)$. Condition $(b)$ holds
    by the definition of $Y_j$ and the induction hypothesis. Finally,
    $Y'\subset A'(D)$ and $F_{2j}$ fixes $Y'$, so $Y'\subset F_{2j}A'(D)=A_j(D)$,
    and $e_j=SGA'(d)=A_j(d)\in A_j(D)$ (in the case $e_j\in Y'$ there is
    nothing to add); this is $(c)$.

    By lemma~\ref{ProductLemma} the operators $A_j$ converge in norm to a
    linear bijection $T$ of $\ell_1$ with
    $\norm{T}\norm{T^{-1}}\leq\prod_m(1+\theta_m){(1-\theta_m)}^{-1}\leq
    1+\varepsilon$. For every $i\in\mathbb{N}$ and $j\geq i$ we have
    $A_j(d_i)\in Y_j$ by $(a)$, and by $(d)$ the operators $F_{2j+1},F_{2j+2},\ldots$
    fix $Y_j$, so $A_k(d_i)=A_j(d_i)$ for all $k\geq j$ and
    $T(d_i)=A_i(d_i)\in Y_i\subset\lo{\mathbb{N}}$. Hence
    $T(D)\subset\lo{\mathbb{N}}$. Similarly, for every $i\in\mathbb{N}$ by
    $(b)$ and $(c)$ there is $d\in D$ with $A_i(d)=e_i\in Y_i$, and by $(d)$
    we get $T(d)=A_i(d)=e_i$. Hence $\lo{\mathbb{N}}\subset T(D)$.
\end{proof}

Now we are ready to prove the main result of the paper.

\begin{theorem}\label{MainTheorem} Let $P$ be a normed space of countable
    (finite or countably infinite) algebraic dimension. Then the following
    conditions are equivalent:
    \begin{enumerate}[label = (\roman*)]
        \item $P$ is extremely projective;
        \item $\overline{P}$ is isometrically isomorphic to $\ellone{\Gamma}$
              for some countable set $\Gamma$;
        \item $\BM(P,\lo{\Gamma})=1$ for some countable set $\Gamma$;
        \item for every $\varepsilon>0$ there is a set $\Lambda$ with
              $\BM(P,\lo{\Lambda})<1+\varepsilon$.
    \end{enumerate}
\end{theorem}
\begin{proof} $(i)\implies(ii)$ By theorem~\ref{CompletionIsEllOne} the space
    $\overline{P}$ is isometrically isomorphic to $\ellone{\Gamma}$ for some
    set $\Gamma$. The rational linear span of a countable Hamel basis of $P$
    is a countable dense subset of $\overline{P}$, so $\overline{P}$ is
    separable and $\Gamma$ is countable.

    $(ii)\implies(iii)$ Let $J:\overline{P}\to\ellone{\Gamma}$ be an isometric
    isomorphism. If $\Gamma$ is finite, then $\overline{P}$ is finite
    dimensional, so $P=\overline{P}$ and $J$ maps $P$ onto
    $\ellone{\Gamma}=\lo{\Gamma}$. If $\Gamma$ is infinite, we may assume that
    $\Gamma=\mathbb{N}$. Then $D=J(P)$ is a dense subspace of $\ell_1$ of
    countable algebraic dimension, and by theorem~\ref{BackAndForth}
    $\BM(P,\lo{\mathbb{N}})=\BM(D,\lo{\mathbb{N}})=1$.

    $(iii)\implies(iv)$ is obvious.

    $(iv)\implies(i)$ Fix $\varepsilon>0$ and a linear bijection
    $S:P\to\lo{\Lambda}$ with $\norm{S}\norm{S^{-1}}<1+\varepsilon$. Then $P$
    is a $(1+\varepsilon)$-retract of $\lo{\Lambda}$ via $S^{-1}$ and $S$, so
    $P$ is extremely projective by theorem~\ref{ExtrProjCriterion}.
\end{proof}

\begin{corollary}\label{DenseCountableSubspaces} Every subspace of $\ell_1$
    of countable algebraic dimension which is dense in $\ell_1$ is extremely
    projective and almost isometric to $\lo{\mathbb{N}}$, that is, has
    Banach--Mazur distance $1$ to $\lo{\mathbb{N}}$.
\end{corollary}

Corollary~\ref{DenseCountableSubspaces} may be compared with the results of
H\'ajek and Russo on densely isomorphic normed spaces
\cite{HajekRussoDenselyIsom}: by \cite[theorem 2.4]{HajekRussoDenselyIsom}
every dense subspace of $\ellone{\Gamma}$ contains a dense subspace
topologically isomorphic to $\lo{\Gamma}$. Corollary~\ref{DenseCountableSubspaces}
says that for dense subspaces of $\ell_1$ of countable algebraic dimension
the subspace itself is almost isometric to $\lo{\mathbb{N}}$.

%-------------------------------------------------------------------------------
%	Counterexample
%-------------------------------------------------------------------------------

\section{Extremely projective spaces which are not free}\label{SectionCounterexample}

In this section we show that extreme projectivity is not an isometric
invariant of the form ``$P$ is isometrically isomorphic to $\lo{\Lambda}$''.
By corollary~\ref{DenseCountableSubspaces}, for every subspace $V$ of $\ell_1$
of countable algebraic dimension the space $\lo{\mathbb{N}}+V$ is extremely
projective. To see that such spaces are not isometrically isomorphic to free normed
spaces whenever $V\not\subset\lo{\mathbb{N}}$, we look at extreme points of
unit balls. By $\ext(B)$ we denote the set of extreme points of a convex set
$B$ and by $\conv(M)$ the convex hull of a set $M$.

\begin{lemma}\label{ExtremePointsLemma} Let $S$ be a set and $E$ be a
    subspace of $\ellone{S}$ containing $\lo{S}$, endowed with the norm of
    $\ellone{S}$. Then
    $\ext(B_E)=\{\lambda e_s:s\in S,\ \lambda\in\mathbb{C},\ |\lambda|=1\}$
    and $\conv(\ext(B_E))=B_{\lo{S}}$.
\end{lemma}
\begin{proof} The points $\lambda e_s$ with $|\lambda|=1$ are extreme points
    of $B_{\ellone{S}}$, and $B_E=B_{\ellone{S}}\cap E$, so they are extreme
    points of $B_E$. Conversely, let $x\in B_E$ be not of this form. If
    $x=0$, then $x=\frac12 e_s+\frac12(-e_s)$ for any $s\in S$. If
    $0<\norm{x}<1$, then $x=\frac12(1+t)x+\frac12(1-t)x$ for $t=1-\norm{x}>0$,
    and both $(1\pm t)x$ belong to $B_E$. If $\norm{x}=1$, then $x$ has at
    least two non-zero coordinates $x_n$ and $x_m$, $n\neq m$. Put
    $t=\min(|x_n|,|x_m|)>0$ and
    $u=t(\operatorname{sgn}(x_n)e_n-\operatorname{sgn}(x_m)e_m)\in\lo{S}\subset E$.
    Then $x\pm u\in E$ and
    \[
    \norm{x\pm u}=(|x_n|\pm t)+(|x_m|\mp t)+\sum_{s\neq n,m}|x_s|=\norm{x}=1,
    \]
    so $x=\frac12(x+u)+\frac12(x-u)$ with $u\neq 0$. In all three cases $x$ is
    not an extreme point of $B_E$.

    Every finite convex combination of the points $\lambda e_s$ is a finitely
    supported element of $B_{\ellone{S}}$, so $\conv(\ext(B_E))\subset
    B_{\lo{S}}$. Conversely, if $x\in B_{\lo{S}}$, then
    \[
    x=\sum_{s\in\supp(x)}|x_s|\,\operatorname{sgn}(x_s)e_s+(1-\norm{x})\cdot 0,
    \]
    and $0=\frac12 e_s+\frac12(-e_s)$, so $x\in\conv(\ext(B_E))$.
\end{proof}

\begin{theorem}\label{NotIsometricToFree} Let $S$ be a set and $E$ be a
    subspace of $\ellone{S}$ containing $\lo{S}$, endowed with the norm of
    $\ellone{S}$. If $E$ is isometrically isomorphic to $\lo{\Lambda}$ for
    some set $\Lambda$, then $E=\lo{S}$.
\end{theorem}
\begin{proof} Let $T:\lo{\Lambda}\to E$ be an isometric isomorphism. Then
    $T(B_{\lo{\Lambda}})=B_E$ and $T$ maps $\ext(B_{\lo{\Lambda}})$ onto
    $\ext(B_E)$. Applying lemma~\ref{ExtremePointsLemma} to $E$ and to
    $\lo{\Lambda}$ itself we get
    \[
    B_E=T(B_{\lo{\Lambda}})=T(\conv(\ext(B_{\lo{\Lambda}})))
    =\conv(T(\ext(B_{\lo{\Lambda}})))=\conv(\ext(B_E))=B_{\lo{S}}.
    \]
    Hence $E=\lo{S}$.
\end{proof}

\begin{corollary}\label{CounterexampleCorollary} Let $V$ be a subspace of
    $\ell_1$ of countable algebraic dimension with $V\not\subset\lo{\mathbb{N}}$,
    for example $V=\mathbb{C}v$ for $v={(2^{-n})}_{n\in\mathbb{N}}$. Then
    $\lo{\mathbb{N}}+V$ is an extremely projective normed space which is not
    isometrically isomorphic to $\lo{\Lambda}$ for any set $\Lambda$.
\end{corollary}
\begin{proof} The space $\lo{\mathbb{N}}+V$ is dense in $\ell_1$ and has
    countable algebraic dimension, so it is extremely projective by
    corollary~\ref{DenseCountableSubspaces}, and it is not isometrically
    isomorphic to any $\lo{\Lambda}$ by theorem~\ref{NotIsometricToFree}.
\end{proof}

Corollary~\ref{CounterexampleCorollary} answers the question of Helemskii
from~\cite[section 2]{HelExtrProjSeqAlg} and~\cite[section 1]{HelMetrFrQMod}
in the negative and disproves \cite[conjecture 3.1]{GronLiftProbNormSp}. Note
that by theorem~\ref{NotIsometricToFree} every subspace
$E$ with $\lo{\mathbb{N}}\subsetneq E\subset\ell_1$ fails to be isometrically
isomorphic to a free normed space, but by example~\ref{WeightedExample} not
all of these subspaces are extremely projective. Thus, between
$\lo{\mathbb{N}}$ and $\ell_1$ there are both extremely projective and
non-projective normed spaces, and the difference between them is invisible on
the level of completions. The next section provides a criterion which
distinguishes them.

%-------------------------------------------------------------------------------
%	Uncountable dimension
%-------------------------------------------------------------------------------

\section{Extremely projective normed spaces of uncountable
  dimension}\label{SectionUncountable}

Condition $(ii)$ of theorem~\ref{MainTheorem} shows that for normed spaces of
countable dimension extreme projectivity is a property of the completion. By
example~\ref{WeightedExample} this is no longer true for spaces of uncountable
dimension: the dimension restriction of theorem~\ref{SeparableSubspaces} is an
independent necessary condition.

By theorems~\ref{SeparableSubspaces} and~\ref{CompletionIsEllOne}, an
extremely projective normed space is, up to an isometric isomorphism, a
dense subspace of some $\ellone{\Gamma}$ all of whose separable subspaces have
countable algebraic dimension; by the theorem of Gr{\o}nb{\ae}k it is also
topologically isomorphic to $\lo{\Gamma}$. For countable $\Gamma$ the first
two conditions are also sufficient by theorems~\ref{SeparableSubspaces}
and~\ref{MainTheorem}. We shall now show that for uncountable $\Gamma$ they
are not. The obstruction is a
set-theoretic refinement of theorem~\ref{SeparableSubspaces}. By a limit
ordinal we always mean a non-zero ordinal which is not a successor. Recall
that a subset of the first uncountable ordinal $\omega_1$ is called closed unbounded
if it is unbounded in $\omega_1$ and contains the supremum of each of its
countable subsets, and stationary if it intersects every closed unbounded
set. The set of limit ordinals below $\omega_1$ is closed unbounded. We shall
use Fodor's lemma \cite[theorem 8.7]{JechSetTheory}: if $E\subset\omega_1$ is
stationary and $\gamma:E\to\omega_1$ satisfies $\gamma(\alpha)<\alpha$ for all
$\alpha\in E$, then $\gamma$ is constant on some stationary subset of $E$.
For a subset $M$ of a normed space $P$ we denote by $\overline{\Lin}(M)$ the
closure of $\Lin(M)$ in $P$.

\begin{theorem}\label{FodorObstruction} Let $P$ be a topologically
    projective normed space, ${(X_\beta)}_{\beta<\omega_1}$ be a family of
    countable subsets of $P$, $E$ be a stationary set of limit ordinals below
    $\omega_1$, and ${(v_\alpha)}_{\alpha\in E}$ be a family in $P$ such that
    $v_\alpha\in\overline{\Lin}\bigl(\bigcup_{\beta<\alpha}X_\beta\bigr)$ for
    every $\alpha\in E$. Then there is a stationary set $E'\subset E$ such that
    $\Lin\{v_\alpha:\alpha\in E'\}$ has countable algebraic dimension. In
    particular, the family ${(v_\alpha)}_{\alpha\in E}$ is linearly dependent.
\end{theorem}
\begin{proof} By proposition~\ref{CanonicalMorphismCriterion}~$(i)$ there is
    a bounded linear operator $\sigma:P\to\lo{B_P}$ with $\pi_P\sigma=1_P$; it
    is injective. For $\alpha<\omega_1$ put
    \[
    S_\alpha=\bigcup_{\beta<\alpha}\ \bigcup_{x\in X_\beta}\supp(\sigma(x)),
    \]
    a countable subset of $B_P$. Let $\alpha\in E$. Then $\sigma(v_\alpha)$
    lies in the closure in $\ellone{B_P}$ of
    $\sigma\bigl(\Lin\bigcup_{\beta<\alpha}X_\beta\bigr)\subset\lo{S_\alpha}$,
    and since the coordinate functionals are continuous, $\sigma(v_\alpha)$ is
    supported in $S_\alpha$. Thus, the finite set $F_\alpha=\supp(\sigma(v_\alpha))$
    is contained in $S_\alpha=\bigcup_{\beta<\alpha}S_{\beta+1}$, and since
    $\alpha$ is a limit ordinal there is $\gamma(\alpha)<\alpha$ with
    $F_\alpha\subset S_{\gamma(\alpha)}$. By Fodor's lemma there are a
    stationary set $E'\subset E$ and $\gamma_0<\omega_1$ with
    $\gamma(\alpha)=\gamma_0$ for all $\alpha\in E'$. Then
    $\sigma(v_\alpha)\in\lo{S_{\gamma_0}}$ for all $\alpha\in E'$. The space
    $\lo{S_{\gamma_0}}$ has countable algebraic dimension and $\sigma$ is
    injective, so $\Lin\{v_\alpha:\alpha\in E'\}$ has countable algebraic
    dimension. As $E'$ is uncountable, the family ${(v_\alpha)}_{\alpha\in E'}$,
    and a fortiori ${(v_\alpha)}_{\alpha\in E}$, is linearly dependent.
\end{proof}

\begin{corollary}\label{OmegaOneContinuity} Let $P$ be a topologically
    projective normed space and ${(X_\beta)}_{\beta<\omega_1}$ be a family of
    countable subsets of $P$. For $\alpha<\omega_1$ put
    $D_\alpha=\overline{\Lin}\bigl(\bigcup_{\beta<\alpha}X_\beta\bigr)$. Then
    the set of limit ordinals $\alpha<\omega_1$ for which
    $D_\alpha\neq\bigcup_{\alpha'<\alpha}D_{\alpha'}$ is not stationary.
\end{corollary}
\begin{proof} Assume that this set $E$ is stationary and for every
    $\alpha\in E$ choose $v_\alpha\in D_\alpha\setminus\bigcup_{\alpha'<\alpha}D_{\alpha'}$.
    The family ${(v_\alpha)}_{\alpha\in E}$ is linearly independent. Indeed,
    let $\sum_{\alpha\in F}c_\alpha v_\alpha=0$ for a finite set $F\subset E$
    and non-zero scalars $c_\alpha$, and let $\alpha_0=\max F$. Since
    $v_{\alpha_0}\neq 0$, the set $F\setminus\{\alpha_0\}$ is non-empty; put
    $\alpha_1=\max(F\setminus\{\alpha_0\})<\alpha_0$. The spaces $D_\alpha$
    increase with $\alpha$, so
    $v_{\alpha_0}=-c_{\alpha_0}^{-1}\sum_{\alpha\in F\setminus\{\alpha_0\}}c_\alpha v_\alpha
    \in D_{\alpha_1}$, which contradicts the choice of $v_{\alpha_0}$. Now
    theorem~\ref{FodorObstruction} gives a contradiction.
\end{proof}

Theorem~\ref{SeparableSubspaces} says that the closed linear span of
countably many vectors of a topologically projective space has countable
dimension; corollary~\ref{OmegaOneContinuity} says that, in addition, the
closed linear spans of initial segments of an $\omega_1$-sequence of vectors
must behave continuously at almost all limit stages. The following example
shows that the second condition is not implied by the first one, even for
dense subspaces of $\ellone{\Gamma}$.

\begin{example}\label{LadderSpace} For every limit ordinal $\alpha<\omega_1$
    fix a strictly increasing sequence ${(\alpha_n)}_{n\in\mathbb{N}}$ of
    ordinals with $\sup_n\alpha_n=\alpha$ and put
    $v_\alpha=\sum_{n=1}^\infty 2^{-n}e_{\alpha_n}\in\ellone{\omega_1}$. Let
    \[
    P_L=\lo{\omega_1}+\Lin\{v_\alpha:\alpha<\omega_1\text{ is a limit ordinal}\}
    \subset\ellone{\omega_1}
    \]
    with the norm inherited from $\ellone{\omega_1}$. Then
    \begin{enumerate}[label = (\roman*)]
        \item $P_L$ is dense in $\ellone{\omega_1}$;
        \item for every countable ordinal $\delta$ the space
              $P_L\cap\ellone{\delta}=\lo{\delta}+\Lin\{v_\alpha:\alpha\leq\delta
              \text{ is a limit ordinal}\}$ is an extremely projective normed
              space of countable algebraic dimension; in particular, every
              separable subspace of $P_L$ has countable algebraic dimension;
        \item $P_L$ is not topologically projective.
    \end{enumerate}
\end{example}
\begin{proof} $(i)$ is obvious, since $P_L\supset\lo{\omega_1}$.

    First we note that for limit ordinals $\alpha<\beta$ the set
    $\supp(v_\alpha)\cap\supp(v_\beta)$ is finite: $\supp(v_\alpha)\subset\alpha$,
    while only finitely many $\beta_n$ are smaller than $\alpha$, because
    $\sup_n\beta_n=\beta>\alpha$.

    $(ii)$ Let $\delta$ be a countable ordinal and let
    $x=y+\sum_{\alpha\in F}c_\alpha v_\alpha\in P_L\cap\ellone{\delta}$, where
    $y\in\lo{\omega_1}$, $F$ is a finite set of limit ordinals and $c_\alpha\neq 0$
    for $\alpha\in F$. Suppose that $F\not\subset[0,\delta]$ and put
    $\alpha_0=\max F>\delta$. The set $\supp(v_{\alpha_0})\cap[\delta,\alpha_0)$
    is infinite, because $\sup_n\alpha_{0,n}=\alpha_0>\delta$, while the sets
    $\supp(y)$ and $\supp(v_\alpha)\cap\supp(v_{\alpha_0})$, $\alpha\in F\setminus\{\alpha_0\}$,
    are finite. Hence $x$ has infinitely many non-zero coordinates in
    $[\delta,\alpha_0)$, which contradicts $x\in\ellone{\delta}$. Therefore
    $F\subset[0,\delta]$, and then also $\supp(y)\subset\delta$. This proves
    the formula for $P_L\cap\ellone{\delta}$. The space $P_L\cap\ellone{\delta}$
    has countable algebraic dimension, and it is dense in $\ellone{\delta}$,
    because it contains $\lo{\delta}$. By theorem~\ref{MainTheorem} it is
    extremely projective. Finally, if $Z$ is a separable subspace of $P_L$,
    then the supports of the elements of $Z$ are contained in a countable
    subset of $\omega_1$, hence in some countable ordinal $\delta$, so
    $Z\subset P_L\cap\ellone{\delta}$ has countable algebraic dimension.

    $(iii)$ Let $E$ be the set of limit ordinals below $\omega_1$; it is
    closed unbounded, hence stationary. For $\alpha\in E$ we have
    $v_\alpha=\lim_N\sum_{n=1}^N 2^{-n}e_{\alpha_n}\in\overline{\Lin}\{e_\beta:\beta<\alpha\}$.
    The family ${(v_\alpha)}_{\alpha\in E}$ is linearly independent: if
    $\sum_{\alpha\in F}c_\alpha v_\alpha=0$ with $F\subset E$ finite and
    $c_\alpha\neq 0$, then for $\alpha_0=\max F$ and all sufficiently large
    $n$ the coordinate of the left hand side at $\alpha_{0,n}$ equals
    $2^{-n}c_{\alpha_0}\neq 0$, because the supports of the other $v_\alpha$
    meet $\supp(v_{\alpha_0})$ in finite sets. By
    theorem~\ref{FodorObstruction}, applied to $X_\beta=\{e_\beta\}$, the space
    $P_L$ is not topologically projective.
\end{proof}

Ladder systems and Fodor's lemma have been used before in Banach space
theory, for instance in the study of the spaces $C(K)$ of ladder system
compacta \cite{AvilesKoszmiderLadder}; here they enter through the dimension
restrictions of topologically projective normed spaces.

Thus, a dense subspace of $\ellone{\Gamma}$ all of whose separable subspaces
have countable algebraic dimension need not be topologically projective, and
the answer to the question of Helemskii on dense subspaces of
$\ellone{\Lambda}$ is more involved for uncountable $\Lambda$. Note that the
space $P_L$ is ``locally'' extremely projective: all its countably generated
pieces $P_L\cap\ellone{\delta}$ are extremely projective. So extreme
projectivity of a normed space is not determined by its separable subspaces.

We now turn to positive results. The obstruction in example~\ref{LadderSpace}
is the discontinuity of the chain $P_L\cap\ellone{\delta}$ at limit stages.
We shall show that for topologically projective spaces such discontinuities
can always be avoided, and that this leads to a complete description of
extremely projective normed spaces of arbitrary dimension. The proof is a
transfinite version of theorem~\ref{BackAndForth} combined with an induction
on the cardinality of $\Gamma$. The key point of the transfinite construction
is that the perturbations used at different steps are controlled on pairwise
disjoint blocks of coordinates, so that their ``product'' has norm bounded
independently of the length of the construction.

Let $\Gamma$ be a set, $\lambda$ be an ordinal and ${(\Gamma_\delta)}_{\delta\leq\lambda}$
be an increasing family of subsets of $\Gamma$ such that $\Gamma_0=\varnothing$,
$\Gamma_\lambda=\Gamma$ and $\Gamma_\delta=\bigcup_{\delta'<\delta}\Gamma_{\delta'}$
for every limit ordinal $\delta\leq\lambda$. We call such a family a
filtration of $\Gamma$ and put $N_\delta=\Gamma_{\delta+1}\setminus\Gamma_\delta$
for $\delta<\lambda$. Then $\Gamma_\delta=\bigsqcup_{\delta'<\delta}N_{\delta'}$
for all $\delta\leq\lambda$, and $\ellone{\Gamma_\delta}$ is the $\ell_1$-sum of
the spaces $\ellone{N_{\delta'}}$, $\delta'<\delta$. For $x\in\ellone{\Gamma}$
and $A\subset\Gamma$ we denote by $x|_A$ the restriction of $x$ to $A$,
extended by zero; note that $x|_A$ need not belong to a given subspace
$P\ni x$. For a subspace $P\subset\ellone{\Gamma}$ we put
$P_\delta=P\cap\ellone{\Gamma_\delta}$ and
\[
D_\delta=\{x|_{N_\delta}:x\in P_{\delta+1}\}\subset\ellone{N_\delta},
\qquad \delta<\lambda.
\]
Since the kernel of the restriction map $x\mapsto x|_{N_\delta}$ on
$P_{\delta+1}$ is $P_{\delta+1}\cap\ellone{\Gamma_\delta}=P_\delta$, the space
$D_\delta$ is isomorphic as a vector space to $P_{\delta+1}/P_\delta$.

\begin{theorem}\label{TransfiniteBackAndForth} Let $\Gamma$ be a set, $P$ be
    a dense subspace of $\ellone{\Gamma}$, ${(\Gamma_\delta)}_{\delta\leq\lambda}$
    be a filtration of $\Gamma$ and $0<\theta<1/2$. Suppose that
    \begin{enumerate}[label = (\roman*)]
        \item $P_\delta$ is dense in $\ellone{\Gamma_\delta}$ for every
              $\delta\leq\lambda$;
        \item $P_\delta=\bigcup_{\delta'<\delta}P_{\delta'}$ for every limit
              ordinal $\delta\leq\lambda$;
        \item for every $\delta<\lambda$ there is a linear bijection $U_\delta$
              of $\ellone{N_\delta}$ with $U_\delta(D_\delta)=\lo{N_\delta}$,
              $\norm{U_\delta}\leq 1+\theta$ and $\norm{U_\delta^{-1}}\leq 1$.
    \end{enumerate}
    Then there exists a linear bijection $T:\ellone{\Gamma}\to\ellone{\Gamma}$
    with $T(P)=\lo{\Gamma}$, $\norm{T}\leq{(1+\theta)}^2$ and
    $\norm{T^{-1}}\leq{(1-\theta)}^{-1}$.
\end{theorem}

\begin{proof} We shall construct by transfinite recursion, for every
    $\delta<\lambda$, a bounded linear operator
    $V_\delta:\ellone{N_\delta}\to\ellone{\Gamma_\delta}$ with
    $\norm{V_\delta}\leq\theta$. For $\delta\leq\lambda$ and
    $x\in\ellone{\Gamma_\delta}$ put
    \[
    W_\delta(x)=x+\sum_{\delta'<\delta}V_{\delta'}(x|_{N_{\delta'}}),\qquad
    U_{<\delta}(x)=\sum_{\delta'<\delta}U_{\delta'}(x|_{N_{\delta'}}),\qquad
    T_\delta=U_{<\delta}W_\delta.
    \]
    Since the sets $N_{\delta'}$ are pairwise disjoint and
    $\norm{x}=\sum_{\delta'<\delta}\norm{x|_{N_{\delta'}}}$, the series defining
    $W_\delta(x)$ converges absolutely in $\ellone{\Gamma_\delta}$ and
    $\norm{W_\delta-1}\leq\theta$, so $W_\delta$ is a linear bijection of
    $\ellone{\Gamma_\delta}$ with $\norm{W_\delta^{-1}}\leq{(1-\theta)}^{-1}$;
    the operator $U_{<\delta}$ is the $\ell_1$-sum of the operators
    $U_{\delta'}$, so it is a linear bijection of $\ellone{\Gamma_\delta}$ with
    $\norm{U_{<\delta}}\leq 1+\theta$ and $\norm{U_{<\delta}^{-1}}\leq 1$.
    Hence $T_\delta$ is a linear bijection of $\ellone{\Gamma_\delta}$ with
    $\norm{T_\delta}\leq{(1+\theta)}^2$ and $\norm{T_\delta^{-1}}\leq{(1-\theta)}^{-1}$.
    Moreover, for $\delta'<\delta$ the operator $T_\delta$ coincides with
    $T_{\delta'}$ on $\ellone{\Gamma_{\delta'}}$, because for
    $x\in\ellone{\Gamma_{\delta'}}$ we have $x|_{N_{\delta''}}=0$ for
    $\delta''\geq\delta'$. The operators $V_\delta$ will be chosen so that
    \begin{equation}\label{InductiveClaim}
    T_\delta(P_\delta)=\lo{\Gamma_\delta}\quad\text{for all }\delta\leq\lambda.
    \end{equation}
    Then $T=T_\lambda$ is the required operator.

    Let $\delta<\lambda$ and suppose that $V_{\delta'}$ has been constructed
    for all $\delta'<\delta$ in such a way that \eqref{InductiveClaim} holds
    for all ordinals $\leq\delta$ (for $\delta=0$ this is trivial, since
    $\Gamma_0=\varnothing$). We construct $V_\delta$ so that
    \eqref{InductiveClaim} holds for $\delta+1$; note that
    \eqref{InductiveClaim} for a limit ordinal $\mu\leq\lambda$ follows from
    \eqref{InductiveClaim} for all $\delta<\mu$, because by $(ii)$ and the
    compatibility of the operators $T_\delta$ we have
    $T_\mu(P_\mu)=\bigcup_{\delta<\mu}T_\delta(P_\delta)
    =\bigcup_{\delta<\mu}\lo{\Gamma_\delta}=\lo{\Gamma_\mu}$.

    If $N_\delta=\varnothing$, then $P_{\delta+1}=P_\delta$,
    $T_{\delta+1}=T_\delta$, and we put $V_\delta=0$. So let
    $N_\delta\neq\varnothing$, and put $U=U_\delta$ and $y_n=U^{-1}(e_n)$ for
    $n\in N_\delta$. By $(iii)$, ${(y_n)}_{n\in N_\delta}$ is a Hamel basis of
    $D_\delta$ and
    \[
    \Bigl\lVert\sum_{n}c_ny_n\Bigr\rVert\geq\norm{U}^{-1}\Bigl\lVert\sum_n c_ne_n\Bigr\rVert
    \geq{(1+\theta)}^{-1}\sum_n|c_n|
    \]
    for all finitely supported families ${(c_n)}_{n\in N_\delta}$ of scalars.
    For every $n\in N_\delta$ choose $x_n\in P_{\delta+1}$ with
    $x_n|_{N_\delta}=y_n$. Since $P_\delta$ is dense in $\ellone{\Gamma_\delta}$
    by $(i)$, the coset $-x_n|_{\Gamma_\delta}+P_\delta$ is dense in
    $\ellone{\Gamma_\delta}$, and we can choose $z_n\in\ellone{\Gamma_\delta}$ with
    \[
    x_n|_{\Gamma_\delta}+z_n\in P_\delta,\qquad\norm{z_n}\leq\theta{(1+\theta)}^{-2}.
    \]
    Define $V_\delta$ on $D_\delta$ by $V_\delta(\sum_n c_ny_n)=\sum_n c_nW_\delta(z_n)$.
    Then
    \[
    \Bigl\lVert V_\delta\Bigl(\sum_n c_ny_n\Bigr)\Bigr\rVert
    \leq\norm{W_\delta}\,\theta{(1+\theta)}^{-2}\sum_n|c_n|
    \leq\theta{(1+\theta)}^{-1}\sum_n|c_n|
    \leq\theta\Bigl\lVert\sum_n c_ny_n\Bigr\rVert,
    \]
    so $V_\delta$ extends by continuity from the dense subspace $D_\delta$ to
    a bounded linear operator $V_\delta:\ellone{N_\delta}\to\ellone{\Gamma_\delta}$
    with $\norm{V_\delta}\leq\theta$ (recall that $D_\delta$ is dense in
    $\ellone{N_\delta}$ by $(i)$, since the restriction map is a contractive
    surjection of $\ellone{\Gamma_{\delta+1}}$ onto $\ellone{N_\delta}$).

    We claim that
    \begin{equation}\label{KeyInclusion}
    x|_{\Gamma_\delta}+W_\delta^{-1}V_\delta(x|_{N_\delta})\in P_\delta
    \quad\text{for every }x\in P_{\delta+1}.
    \end{equation}
    Indeed, write $x|_{N_\delta}=\sum_n c_ny_n$ with finitely many non-zero
    $c_n$. Then $x-\sum_n c_nx_n\in P_{\delta+1}$ vanishes on $N_\delta$, so it
    belongs to $P_\delta$, and hence
    $x|_{\Gamma_\delta}=\sum_n c_nx_n|_{\Gamma_\delta}+p$ for some $p\in P_\delta$.
    Therefore
    \[
    x|_{\Gamma_\delta}+W_\delta^{-1}V_\delta(x|_{N_\delta})
    =\sum_n c_n\bigl(x_n|_{\Gamma_\delta}+z_n\bigr)+p\in P_\delta.
    \]

    Now let $x\in P_{\delta+1}$. Since $W_{\delta+1}(x)=W_\delta(x|_{\Gamma_\delta})
    +V_\delta(x|_{N_\delta})+x|_{N_\delta}$ and $U_{<\delta+1}$ acts as
    $U_{<\delta}$ on $\ellone{\Gamma_\delta}$ and as $U$ on
    $\ellone{N_\delta}$, we get
    \[
    T_{\delta+1}(x)=U_{<\delta}W_\delta\bigl(x|_{\Gamma_\delta}+W_\delta^{-1}V_\delta(x|_{N_\delta})\bigr)
    +U(x|_{N_\delta})
    =T_\delta\bigl(x|_{\Gamma_\delta}+W_\delta^{-1}V_\delta(x|_{N_\delta})\bigr)
    +U(x|_{N_\delta}).
    \]
    By \eqref{KeyInclusion} and \eqref{InductiveClaim} for $\delta$, the
    first summand belongs to $\lo{\Gamma_\delta}$, and the second one belongs
    to $U(D_\delta)=\lo{N_\delta}$. Hence $T_{\delta+1}(P_{\delta+1})\subset\lo{\Gamma_{\delta+1}}$.
    Conversely, $\lo{\Gamma_\delta}=T_\delta(P_\delta)=T_{\delta+1}(P_\delta)
    \subset T_{\delta+1}(P_{\delta+1})$, and for every $n\in N_\delta$ the
    above formula gives $T_{\delta+1}(x_n)\in\lo{\Gamma_\delta}+U(y_n)
    =\lo{\Gamma_\delta}+e_n$, so $e_n\in T_{\delta+1}(P_{\delta+1})$. Thus,
    $T_{\delta+1}(P_{\delta+1})=\lo{\Gamma_{\delta+1}}$, which completes the
    recursion and the proof.
\end{proof}

If all the sets $N_\delta$ are at most countable and $P_{\delta+1}/P_\delta$
has at most countable algebraic dimension for every $\delta<\lambda$, then
condition $(iii)$ is automatically satisfied for every $\theta>0$: indeed,
$D_\delta$ is a dense subspace of $\ellone{N_\delta}$ of countable algebraic
dimension, so for infinite $N_\delta$ theorem~\ref{BackAndForth} provides a
linear bijection $U_\delta$ with $U_\delta(D_\delta)=\lo{N_\delta}$ and
$\norm{U_\delta}\norm{U_\delta^{-1}}\leq 1+\theta$, which can be normalized by
a scalar, and for finite $N_\delta$ we have $D_\delta=\ellone{N_\delta}$ and
$U_\delta=1_{\ellone{N_\delta}}$ works. Example~\ref{LadderSpace} shows that
condition $(ii)$ cannot be dropped: for the filtration $\Gamma_\delta=\delta$
of $\omega_1$ the space $P_L$ satisfies $(i)$ and $(iii)$, but fails $(ii)$
at every countable limit ordinal (at $\delta=\omega_1$ itself condition
$(ii)$ holds trivially, since every element of $\ellone{\omega_1}$ has
countable support).

We proceed to the construction of filtrations satisfying the hypotheses of
theorem~\ref{TransfiniteBackAndForth}. For the rest of this section we fix a
set $\Gamma$, a dense topologically
projective subspace $P$ of $\ellone{\Gamma}$, a set $S$ and bounded linear
operators $\sigma:P\to\lo{S}$ and $\pi:\lo{S}\to P$ with $\pi\sigma=1_P$ and
$\norm{\pi}\leq 1$; these exist by proposition~\ref{CanonicalMorphismCriterion}.
Since $P$ is dense in $\ellone{\Gamma}$, we can also fix for every
$\gamma\in\Gamma$ a sequence ${(p^\gamma_n)}_{n\in\mathbb{N}}$ in $P$ converging
to $e_\gamma$. For $A\subset\Gamma$ we put $P_A=P\cap\ellone{A}$.

\begin{definition}\label{ClosedPair} A pair $(A,B)$ of subsets $A\subset\Gamma$
    and $B\subset S$ is called closed if
    \begin{enumerate}[label = (\alph*)]
        \item $\supp(\pi(e_s))\subset A$ for all $s\in B$;
        \item $\supp(\sigma(x))\subset B$ for all $x\in P_A$;
        \item $\supp(p^\gamma_n)\subset A$ for all $\gamma\in A$ and
              $n\in\mathbb{N}$.
    \end{enumerate}
\end{definition}

\begin{lemma}\label{ClosedPairProps} Let $(A,B)$ be a closed pair. Then
    $P_A=\pi(\lo{B})$, the space $P_A$ is dense in $\ellone{A}$, and
    $\sigma|_{P_A}:P_A\to\lo{B}$, $\pi|_{\lo{B}}:\lo{B}\to P_A$ exhibit $P_A$ as
    a retract of $\lo{B}$. In particular, $P_A$ is topologically projective.
\end{lemma}
\begin{proof} By $(a)$, $\pi(\lo{B})\subset P\cap\ellone{A}=P_A$. If $x\in P_A$,
    then $\sigma(x)\in\lo{B}$ by $(b)$ and $x=\pi\sigma(x)\in\pi(\lo{B})$. By
    $(c)$, $P_A$ contains the sequences $p^\gamma_n$, $\gamma\in A$, so its
    closure contains $\lo{A}$ and $P_A$ is dense in $\ellone{A}$. The last
    claim is now obvious.
\end{proof}

\begin{lemma}\label{ClosingOff} Let $\mu$ be an infinite cardinal,
    $A\subset\Gamma$ and $B\subset S$ be sets of cardinality at most $\mu$.
    Then there is a closed pair $(A',B')$ with $A\subset A'$, $B\subset B'$ and
    $|A'|,|B'|\leq\mu$.
\end{lemma}
\begin{proof} Define increasing sequences $A=A_0\subset A_1\subset\ldots$ and
    $B=B_0\subset B_1\subset\ldots$ as follows. Given $A_k$ and $B_k$ of
    cardinality at most $\mu$, let $H_k$ be a Hamel basis of $P_{A_k}$. Since
    $P_{A_k}\subset\ellone{A_k}$ has density character at most $\mu$ (the
    density character of a metric space is inherited by its subspaces), by
    theorem~\ref{SeparableSubspaces} $|H_k|\leq\mu$. Put
    \[
    B_{k+1}=B_k\cup\bigcup_{x\in H_k}\supp(\sigma(x)),\qquad
    A_{k+1}=A_k\cup\bigcup_{s\in B_{k+1}}\supp(\pi(e_s))
    \cup\bigcup_{\gamma\in A_k,\,n\in\mathbb{N}}\supp(p^\gamma_n).
    \]
    All supports are countable, so $|A_{k+1}|,|B_{k+1}|\leq\mu$. Put
    $A'=\bigcup_k A_k$ and $B'=\bigcup_k B_k$. Conditions $(a)$ and $(c)$ hold
    for $(A',B')$, since every $s\in B'$ and every $\gamma\in A'$ belongs to
    some $B_{k+1}$ or $A_k$ respectively. To verify $(b)$, note first that
    $\bigcup_k P_{A_k}$ is dense in $\ellone{A'}$: indeed, $P_{A_{k+1}}$
    contains $p^\gamma_n$ for all $\gamma\in A_k$, so the closure of
    $\bigcup_k P_{A_k}$ contains $\lo{A'}$. Let $x\in P_{A'}$ and choose
    $x_j\in\bigcup_k P_{A_k}$ with $x_j\to x$. Then $\sigma(x_j)\in\lo{B_{k_j+1}}
    \subset\ellone{B'}$ for suitable $k_j$, because $\supp(\sigma(y))\subset
    B_{k+1}$ for $y\in P_{A_k}$ by the choice of $B_{k+1}$. Since $\sigma$ is
    continuous and $\ellone{B'}$ is closed in $\ellone{S}$, we get
    $\sigma(x)\in\ellone{B'}\cap\lo{S}=\lo{B'}$.
\end{proof}

Note that in the last argument the support of $x\in P_{A'}$ need not be
contained in a single $A_k$; this is why density was used instead of a
direct argument. The same argument proves the following lemma.

\begin{lemma}\label{UnionsOfClosedPairs} Let ${((A_\delta,B_\delta))}_{\delta<\lambda}$
    be an increasing family of closed pairs indexed by a limit ordinal
    $\lambda$, and let $A_\lambda=\bigcup_{\delta<\lambda}A_\delta$,
    $B_\lambda=\bigcup_{\delta<\lambda}B_\delta$. Then $(A_\lambda,B_\lambda)$
    is a closed pair and $P_{A_\lambda}=\bigcup_{\delta<\lambda}P_{A_\delta}$.
\end{lemma}
\begin{proof} Conditions $(a)$ and $(c)$ are inherited from the pairs
    $(A_\delta,B_\delta)$. By lemma~\ref{ClosedPairProps} each $P_{A_\delta}$
    is dense in $\ellone{A_\delta}$, so $\bigcup_{\delta<\lambda}P_{A_\delta}$
    is dense in $\ellone{A_\lambda}$. If $x\in P_{A_\lambda}$, choose
    $x_j\in\bigcup_\delta P_{A_\delta}$ with $x_j\to x$; then
    $\sigma(x_j)\in\ellone{B_\lambda}$ by $(b)$ for the pairs
    $(A_\delta,B_\delta)$, hence $\sigma(x)\in\ellone{B_\lambda}\cap\lo{S}=\lo{B_\lambda}$.
    This is $(b)$ for $(A_\lambda,B_\lambda)$. Finally, by
    lemma~\ref{ClosedPairProps},
    $P_{A_\lambda}=\pi(\lo{B_\lambda})=\bigcup_{\delta<\lambda}\pi(\lo{B_\delta})
    =\bigcup_{\delta<\lambda}P_{A_\delta}$.
\end{proof}

\begin{lemma}\label{BlockRetract} Let $(A,B)\subset(A',B')$ be closed pairs,
    that is, $A\subset A'$ and $B\subset B'$, and let $N=A'\setminus A$. Then
    $D=\{x|_N:x\in P_{A'}\}$ is a dense subspace of $\ellone{N}$ which is a
    retract of $\lo{B'\setminus B}$. In particular, $D$ is topologically
    projective.
\end{lemma}
\begin{proof} The restriction map $r:\ellone{A'}\to\ellone{N}$, $x\mapsto x|_N$,
    is a contractive surjection and $P_{A'}$ is dense in $\ellone{A'}$ by
    lemma~\ref{ClosedPairProps}, so $D=r(P_{A'})$ is dense in $\ellone{N}$.
    The kernel of $r|_{P_{A'}}$ is $P_{A'}\cap\ellone{A}=P_A$.

    Let $y\in D$ and $x\in P_{A'}$ with $x|_N=y$. By $(b)$ for $(A',B')$ we
    have $\sigma(x)\in\lo{B'}$. If $x'\in P_{A'}$ is another element with
    $x'|_N=y$, then $x-x'\in P_A$, so $\sigma(x-x')\in\lo{B}$ by $(b)$ for
    $(A,B)$, and $\sigma(x)|_{S\setminus B}=\sigma(x')|_{S\setminus B}$.
    Therefore the map
    \[
    \rho:D\to\lo{B'\setminus B},\qquad\rho(y)=\sigma(x)|_{S\setminus B}
    \quad(x\in P_{A'},\ x|_N=y),
    \]
    is well defined; it is linear, because if $x_1,x_2\in P_{A'}$ restrict
    to $y_1,y_2$, then $x_1+cx_2$ restricts to $y_1+cy_2$. It is bounded:
    given $y\in D$, $x_0\in P_{A'}$
    with $x_0|_N=y$ and $\varepsilon>0$, use the density of $P_A$ in
    $\ellone{A}$ to find $p\in P_A$ with $\norm{x_0|_A+p}\leq\varepsilon$;
    then $x=x_0+p\in P_{A'}$ satisfies $x|_N=y$ and
    $\norm{x}=\norm{x|_A}+\norm{y}\leq\norm{y}+\varepsilon$, so
    $\norm{\rho(y)}\leq\norm{\sigma(x)}\leq\norm{\sigma}(\norm{y}+\varepsilon)$.
    Hence $\norm{\rho}\leq\norm{\sigma}$. Finally, let $r\pi:\lo{B'\setminus B}\to\ellone{N}$
    be the composition of $\pi$ with the restriction to $N$. For $y=x|_N$ as
    above,
    \[
    r\pi\rho(y)=r\bigl(\pi\sigma(x)-\pi(\sigma(x)|_B)\bigr)
    =r(x)-r\bigl(\pi(\sigma(x)|_B)\bigr)=y,
    \]
    because $\pi(\sigma(x)|_B)\in\pi(\lo{B})=P_A\subset\ellone{A}$ vanishes on
    $N$. Thus, $r\pi$ maps $\lo{B'\setminus B}$ into $D$ (its range is
    contained in $r(P_{A'})=D$) and $r\pi\rho=1_D$; so $D$ is a retract of
    $\lo{B'\setminus B}$ and is topologically projective by
    proposition~\ref{RetractOfProj}.
\end{proof}

\begin{theorem}\label{GeneralTheorem} Let $\Gamma$ be a set and $P$ be a
    dense topologically projective subspace of $\ellone{\Gamma}$. Then for
    every $\theta>0$ there is a linear bijection $T$ of $\ellone{\Gamma}$ with
    $T(P)=\lo{\Gamma}$, $\norm{T}\leq 1+\theta$ and $\norm{T^{-1}}\leq 1$. In
    particular, $\BM(P,\lo{\Gamma})=1$ and $P$ is extremely projective.
\end{theorem}
\begin{proof} We proceed by induction on the cardinality $\kappa$ of
    $\Gamma$. If $\Gamma$ is finite, then $P=\ellone{\Gamma}=\lo{\Gamma}$ and
    $T=1_{\ellone{\Gamma}}$ works. If $\Gamma$ is countably infinite, then
    $P$ has countable algebraic dimension by theorem~\ref{SeparableSubspaces},
    and theorem~\ref{BackAndForth} provides a linear bijection $T$ with
    $T(P)=\lo{\Gamma}$ and $\norm{T}\norm{T^{-1}}\leq 1+\theta$; multiplying
    $T$ by a scalar we get $\norm{T^{-1}}\leq 1$ and $\norm{T}\leq 1+\theta$.

    Let $\kappa$ be uncountable and assume that the theorem holds for all
    sets of cardinality less than $\kappa$. Since the conclusion for a smaller
    $\theta$ is stronger, we may fix $\theta\in(0,1/2)$; put
    $\theta'=\theta/8$, so that ${(1+\theta')}^2{(1-\theta')}^{-1}\leq 1+\theta$.
    Fix $S$, $\sigma$, $\pi$ and the sequences $p^\gamma_n$ as above. Let
    $\lambda$ be the cofinality of $\kappa$ and ${(\Gamma^\delta)}_{\delta<\lambda}$
    be a family of subsets of $\Gamma$ of cardinality less than $\kappa$ with
    $\bigcup_{\delta<\lambda}\Gamma^\delta=\Gamma$. Such a family exists: if
    $\kappa$ is regular, then $\lambda=\kappa$ and we may take
    $\Gamma^\delta=\{\gamma_\delta\}$ for an enumeration
    ${(\gamma_\delta)}_{\delta<\kappa}$ of $\Gamma$; if $\kappa$ is singular,
    we take an enumeration ${(\gamma_\alpha)}_{\alpha<\kappa}$ of $\Gamma$, an
    increasing sequence of cardinals ${(\kappa_\delta)}_{\delta<\lambda}$ with
    supremum $\kappa$, and $\Gamma^\delta=\{\gamma_\alpha:\alpha<\kappa_\delta\}$.
    By transfinite recursion
    we define an increasing family ${((A_\delta,B_\delta))}_{\delta\leq\lambda}$
    of closed pairs with $|A_\delta|,|B_\delta|<\kappa$ for $\delta<\lambda$:
    put $A_0=B_0=\varnothing$; given $(A_\delta,B_\delta)$ with $\delta<\lambda$,
    let $(A_{\delta+1},B_{\delta+1})$ be a closed pair containing
    $(A_\delta\cup\Gamma^\delta,B_\delta)$ with
    $|A_{\delta+1}|,|B_{\delta+1}|\leq\max(|A_\delta|,|B_\delta|,|\Gamma^\delta|,\aleph_0)<\kappa$,
    which exists by lemma~\ref{ClosingOff}; for a limit ordinal
    $\delta\leq\lambda$ let $A_\delta$ and $B_\delta$ be the unions of the
    previous sets. For a limit ordinal $\delta<\lambda$ the set $A_\delta$ is
    a union of fewer than $\operatorname{cf}(\kappa)$ sets of cardinality
    less than $\kappa$, so $|A_\delta|<\kappa$, and similarly $|B_\delta|<\kappa$;
    by lemma~\ref{UnionsOfClosedPairs} the pairs $(A_\delta,B_\delta)$ are
    closed for all $\delta\leq\lambda$. Finally, $A_\lambda=\Gamma$.

    Thus, ${(A_\delta)}_{\delta\leq\lambda}$ is a filtration of $\Gamma$, and
    we verify the hypotheses of theorem~\ref{TransfiniteBackAndForth} for
    this filtration and $\theta'$, where $P_\delta=P_{A_\delta}$ and
    $N_\delta=A_{\delta+1}\setminus A_\delta$. Condition $(i)$ holds by
    lemma~\ref{ClosedPairProps}, and condition $(ii)$ holds by
    lemma~\ref{UnionsOfClosedPairs}. For $(iii)$, let $\delta<\lambda$. By
    lemma~\ref{BlockRetract} the space $D_\delta=\{x|_{N_\delta}:x\in P_{\delta+1}\}$
    is a dense topologically projective subspace of $\ellone{N_\delta}$, and
    $|N_\delta|\leq|A_{\delta+1}|<\kappa$. By the induction hypothesis there
    is a linear bijection $U_\delta$ of $\ellone{N_\delta}$ with
    $U_\delta(D_\delta)=\lo{N_\delta}$, $\norm{U_\delta}\leq 1+\theta'$ and
    $\norm{U_\delta^{-1}}\leq 1$ (if $N_\delta=\varnothing$ there is nothing to
    do). Now theorem~\ref{TransfiniteBackAndForth} gives a linear bijection
    $T$ of $\ellone{\Gamma}$ with $T(P)=\lo{\Gamma}$ and
    $\norm{T}\norm{T^{-1}}\leq{(1+\theta')}^2{(1-\theta')}^{-1}\leq 1+\theta$;
    after multiplication by a scalar, $\norm{T^{-1}}\leq 1$ and
    $\norm{T}\leq 1+\theta$. The last claim of the theorem follows from
    theorem~\ref{ExtrProjCriterion}.
\end{proof}

Note that the filtration constructed in the proof of
theorem~\ref{GeneralTheorem} has blocks $N_\delta$ of cardinality less than
$|\Gamma|$, not necessarily countable ones; this is why
theorem~\ref{TransfiniteBackAndForth} was formulated for arbitrary blocks.

Together with theorem~\ref{CompletionIsEllOne}, theorem~\ref{GeneralTheorem}
gives the main result of the paper.

\begin{theorem}\label{FinalCharacterization} For a normed space $P$ the
    following conditions are equivalent:
    \begin{enumerate}[label = (\roman*)]
        \item $P$ is extremely projective;
        \item $P$ is topologically projective and $\overline{P}$ is
              isometrically isomorphic to $\ellone{\Gamma}$ for some set
              $\Gamma$;
        \item $\BM(P,\lo{\Gamma})=1$ for some set $\Gamma$;
        \item for every $\varepsilon>0$ there is a set $\Lambda$ with
              $\BM(P,\lo{\Lambda})<1+\varepsilon$.
    \end{enumerate}
    Moreover, in $(iii)$ the set $\Gamma$ necessarily has the cardinality of
    the algebraic dimension of $P$, and the Banach--Mazur distance in
    $(iii)$ is realized by linear bijections
    $T:\overline{P}\to\ellone{\Gamma}$ with $T(P)=\lo{\Gamma}$ and
    $\norm{T}\norm{T^{-1}}$ arbitrarily close to $1$.
\end{theorem}
\begin{proof} $(i)\implies(ii)$ Extremely projective spaces are
    topologically projective, and their completions are described by
    theorem~\ref{CompletionIsEllOne}.

    $(ii)\implies(iii)$ Let $J:\overline{P}\to\ellone{\Gamma}$ be an isometric
    isomorphism. Then $J(P)$ is a dense topologically projective subspace of
    $\ellone{\Gamma}$, and theorem~\ref{GeneralTheorem} gives
    $\BM(P,\lo{\Gamma})=\BM(J(P),\lo{\Gamma})=1$, realized by linear
    bijections of the required form.

    $(iii)\implies(iv)$ is obvious, and $(iv)\implies(i)$ follows from
    theorem~\ref{ExtrProjCriterion}, since a linear bijection $S:P\to\lo{\Lambda}$
    with $\norm{S}\norm{S^{-1}}<1+\varepsilon$ makes $P$ a
    $(1+\varepsilon)$-retract of $\lo{\Lambda}$. Finally, if
    $\BM(P,\lo{\Gamma})<\infty$, then $P$ and $\lo{\Gamma}$ are isomorphic as
    vector spaces, and the algebraic dimension of $\lo{\Gamma}$ is the
    cardinality of $\Gamma$; so the set $\Gamma$ in $(iii)$ is determined up
    to a bijection, and the realization of the distance obtained in
    $(ii)\implies(iii)$ applies to it.
\end{proof}

\begin{corollary}\label{GronbaekConjecture} Every extremely projective
    normed space $P$ is almost isometric to $\lo{\Gamma}$, where $\Gamma$ is a
    set whose cardinality is the algebraic dimension of $P$. In particular,
    the conjecture of Gr{\o}nb{\ae}k \cite[conjecture 3.1]{GronLiftProbNormSp}
    holds in the almost isometric form, while corollary~\ref{CounterexampleCorollary}
    shows that it fails in the isometric form.
\end{corollary}

We finish with some comments. Let us compare our argument with the proof
of the theorem of Gr{\o}nb{\ae}k, which consists of two steps: a transfinite
construction which finds in every subspace $Y$ of $\lo{\Lambda}$ a
complemented subspace isomorphic to $\lo{\Lambda'}$, where the cardinality of
$\Lambda'$ is the density character of $Y$
\cite[proposition 2.2]{GronLiftProbNormSp}, and the decomposition method of
Pe{\l}czy\'nski \cite[proposition 2.3]{GronLiftProbNormSp}. An inspection of
the first step shows that it produces, for every $\varepsilon>0$, a subspace
which is $(1+\varepsilon)$-isomorphic to $\lo{\Lambda'}$ and
$(1+\varepsilon)$-complemented, so the constants are lost only in the second
step. Theorems~\ref{BackAndForth} and~\ref{TransfiniteBackAndForth} replace
the decomposition method by back-and-forth constructions which keep the
constants, at the price of the continuity condition $(ii)$, and
lemmas~\ref{ClosingOff} and~\ref{UnionsOfClosedPairs} show that this
condition can always be achieved for topologically projective spaces. Next,
theorem~\ref{FinalCharacterization}
shows that for dense subspaces of the spaces $\ellone{\Gamma}$ topological and
extreme projectivity coincide. This is a feature of the exact norm of
$\ellone{\Gamma}$: if $E$ is an infinite dimensional Banach space which is
topologically but not almost isometrically isomorphic to $\ellone{\Gamma}$,
for instance $\ellone{\Gamma}$ with a suitable equivalent norm, then every
dense subspace of $E$ of the form $J(\lo{\Gamma})$, $J$ a topological
isomorphism, is topologically but not extremely projective by
theorem~\ref{CompletionIsEllOne}. Further, by theorem~\ref{FinalCharacterization}
and example~\ref{LadderSpace}, the question of Helemskii on dense subspaces of
$\ellone{\Lambda}$ has the following answer: a dense subspace $P$ of
$\ellone{\Lambda}$ is extremely projective if and only if it is
topologically projective, if and only if it admits a filtration with the
properties $(i)$--$(iii)$ of theorem~\ref{TransfiniteBackAndForth}; for
countable $\Lambda$ this means that $P$ has countable algebraic dimension, but
for uncountable $\Lambda$ the continuity condition $(ii)$ is a genuine
restriction, as the ladder space shows. The positive results of this section
use only the closing-off arguments of lemmas~\ref{ClosingOff}
and~\ref{UnionsOfClosedPairs}. In the next section we sharpen the last
answer: condition $(iii)$ of theorem~\ref{TransfiniteBackAndForth} can be
replaced by a dimension count, and for $|\Lambda|=\aleph_1$ the answer takes
the form of a criterion in terms of stationary sets.

%-------------------------------------------------------------------------------
%	Filtrations with countable blocks
%-------------------------------------------------------------------------------

\section{Filtrations with countable blocks and the case of dimension
  $\aleph_1$}\label{SectionIntrinsic}

In this section we return to the question of Helemskii on dense subspaces of
$\ellone{\Gamma}$ and describe the topologically projective ones in terms of
the coordinate structure of $\ellone{\Gamma}$ only, without reference to a
retraction onto a free space. Note that this structure is intrinsic: by
lemma~\ref{ExtremePointsLemma}, applied to $E=\ellone{\Gamma}$, the extreme
points of the unit ball of $\ellone{\Gamma}$ are exactly the vectors
$\lambda e_\gamma$ with $|\lambda|=1$, so the family of the one dimensional
subspaces $\mathbb{C}e_\gamma$, $\gamma\in\Gamma$, is determined by the norm
of $\ellone{\Gamma}$. As before, for a subspace $P$ of $\ellone{\Gamma}$ and
$A\subset\Gamma$ we write $P_A=P\cap\ellone{A}$, and for a filtration
${(\Gamma_\delta)}_{\delta\leq\lambda}$ of $\Gamma$ we write
$P_\delta=P_{\Gamma_\delta}$.

We begin with a reformulation of theorem~\ref{SeparableSubspaces}. Its proof
uses only the existence of a bounded injective operator into a free normed
space, so it gives the following statement.

\begin{lemma}\label{InjectionDimension} Let $Z$ be a normed space, $S$ be a
    set and $\rho:Z\to\lo{S}$ be an injective bounded linear operator. Then
    the algebraic dimension of $Z$ is at most $\max(\operatorname{dens}(Z),\aleph_0)$.
\end{lemma}
\begin{proof} Let $M$ be a dense subset of $Z$ of cardinality
    $\operatorname{dens}(Z)$ and $S_0=\bigcup_{x\in M}\supp(\rho(x))$, a set
    of cardinality at most $\max(\operatorname{dens}(Z),\aleph_0)$. For
    $z\in Z$ the vector $\rho(z)$ is the limit in $\ellone{S}$ of a sequence
    of elements of $\rho(M)$, which are supported in $S_0$, so $\rho(z)$ is
    supported in $S_0$ by the continuity of the coordinate functionals. Thus,
    $\rho$ is an injective linear map of $Z$ into $\lo{S_0}$, whose algebraic
    dimension is the cardinality of $S_0$.
\end{proof}

For the next two lemmas we return to the standing assumptions of
section~\ref{SectionUncountable}: $P$ is a dense topologically projective
subspace of $\ellone{\Gamma}$, the operators $\sigma:P\to\lo{S}$ and
$\pi:\lo{S}\to P$ satisfy $\pi\sigma=1_P$ and $\norm{\pi}\leq 1$, the
sequences ${(p^\gamma_n)}_{n\in\mathbb{N}}$ in $P$ converge to $e_\gamma$,
and closed pairs are understood in the sense of definition~\ref{ClosedPair}.

\begin{lemma}\label{CountableQuotient} Let $(A,B)$ be a closed pair and
    $A'\subset\Gamma$ be a set with $A\subset A'$ such that $A'\setminus A$
    is at most countable. Then the algebraic dimension of $P_{A'}/P_A$ is at
    most countable.
\end{lemma}
\begin{proof} Put $N=A'\setminus A$ and let $r:\ellone{A'}\to\ellone{N}$ be
    the restriction map. The kernel of $r|_{P_{A'}}$ is $P_{A'}\cap\ellone{A}=P_A$,
    so $D=r(P_{A'})$ is isomorphic as a vector space to $P_{A'}/P_A$, and $D$
    is separable as a subspace of $\ellone{N}$. As in the proof of
    lemma~\ref{BlockRetract}, the map
    \[
    \rho:D\to\lo{S\setminus B},\qquad\rho(y)=\sigma(x)|_{S\setminus B}
    \quad(x\in P_{A'},\ x|_N=y),
    \]
    is well defined, linear and bounded. Indeed, two elements $x,x'\in P_{A'}$
    with $x|_N=x'|_N$ differ by an element of $P_A$, whose image under
    $\sigma$ is supported in $B$ by condition $(b)$ for $(A,B)$, so $\rho$ is
    well defined and linear; and given $y\in D$, $x_0\in P_{A'}$ with
    $x_0|_N=y$ and $\varepsilon>0$, the density of $P_A$ in $\ellone{A}$
    (lemma~\ref{ClosedPairProps}) gives $p\in P_A$ with
    $\norm{x_0|_A+p}\leq\varepsilon$, and then $x=x_0+p\in P_{A'}$ satisfies
    $x|_N=y$ and $\norm{x}=\norm{x_0|_A+p}+\norm{y}\leq\norm{y}+\varepsilon$,
    whence $\norm{\rho(y)}\leq\norm{\sigma}(\norm{y}+\varepsilon)$. The map
    $\rho$ is injective: if $\rho(y)=0$ for $y=x|_N$, then
    $\sigma(x)\in\lo{B}$, so $x=\pi\sigma(x)\in\pi(\lo{B})=P_A$ by
    lemma~\ref{ClosedPairProps}, and $y=x|_N=0$. By
    lemma~\ref{InjectionDimension} the algebraic dimension of $D$ is at most
    countable.
\end{proof}

\begin{lemma}\label{CountableClosingOff} Let $(A,B)$ be a closed pair and
    $N\subset\Gamma$ be an at most countable set. Then there is a closed pair
    $(A',B')$ with $A\cup N\subset A'$, $B\subset B'$ such that
    $A'\setminus A$ and $B'\setminus B$ are at most countable.
\end{lemma}
\begin{proof} We define increasing sequences $A\cup N=A_0\subset A_1\subset\ldots$
    and $B=B_0\subset B_1\subset\ldots$ such that $A_k\setminus A$ and
    $B_k\setminus B$ are at most countable. Given $A_k$ and $B_k$, by
    lemma~\ref{CountableQuotient} there is an at most countable set
    $H_k\subset P_{A_k}$ such that $P_{A_k}=P_A+\Lin(H_k)$. Put
    \[
    B_{k+1}=B_k\cup\bigcup_{x\in H_k}\supp(\sigma(x)),\qquad
    A_{k+1}=A_k\cup\bigcup_{s\in B_{k+1}\setminus B}\supp(\pi(e_s))
    \cup\bigcup_{\gamma\in A_k\setminus A,\,n\in\mathbb{N}}\supp(p^\gamma_n).
    \]
    The sets $B_{k+1}\setminus B$ and $A_{k+1}\setminus A$ are at most
    countable, since all supports involved are at most countable. By
    condition $(b)$ for $(A,B)$ we have $\supp(\sigma(x))\subset B_{k+1}$ for
    all $x\in P_{A_k}$. Put $A'=\bigcup_k A_k$ and $B'=\bigcup_k B_k$. Then
    $A'\setminus A$ and $B'\setminus B$ are at most countable. Condition
    $(a)$ holds for $(A',B')$: for $s\in B$ we have $\supp(\pi(e_s))\subset A$
    by $(a)$ for $(A,B)$, and for $s\in B_{k+1}\setminus B$ we have
    $\supp(\pi(e_s))\subset A_{k+1}$. Similarly, condition $(c)$ holds for
    $(A',B')$, since $\supp(p^\gamma_n)\subset A$ for $\gamma\in A$ and
    $\supp(p^\gamma_n)\subset A_{k+1}$ for $\gamma\in A_k\setminus A$.
    Finally, $\bigcup_k P_{A_k}$ is dense in $\ellone{A'}$: its closure
    contains $P_A$, which is dense in $\ellone{A}$ by
    lemma~\ref{ClosedPairProps}, and the vectors $p^\gamma_n\in P_{A_{k+1}}$,
    $\gamma\in A_k\setminus A$, so it contains $e_\gamma$ for all
    $\gamma\in A'$. If $x\in P_{A'}$, choose $x_j\in\bigcup_k P_{A_k}$ with
    $x_j\to x$; then $\sigma(x_j)\in\lo{B_{k_j+1}}\subset\ellone{B'}$ for
    suitable $k_j$, and since $\sigma$ is continuous and $\ellone{B'}$ is
    closed in $\ellone{S}$, we get $\sigma(x)\in\ellone{B'}\cap\lo{S}=\lo{B'}$.
    This is condition $(b)$ for $(A',B')$.
\end{proof}

\begin{theorem}\label{CountableBlocksCriterion} Let $\Gamma$ be a set and
    $P$ be a dense subspace of $\ellone{\Gamma}$. Then the following
    conditions are equivalent:
    \begin{enumerate}[label = (\roman*)]
        \item $P$ is extremely projective;
        \item $P$ is topologically projective;
        \item there are an ordinal $\lambda$ and a filtration
              ${(\Gamma_\delta)}_{\delta\leq\lambda}$ of $\Gamma$, all of whose
              blocks $N_\delta=\Gamma_{\delta+1}\setminus\Gamma_\delta$ are at
              most countable, such that
              \begin{enumerate}[label = (\alph*)]
                  \item $P_\delta$ is dense in $\ellone{\Gamma_\delta}$ for
                        every $\delta\leq\lambda$;
                  \item $P_\delta=\bigcup_{\delta'<\delta}P_{\delta'}$ for
                        every limit ordinal $\delta\leq\lambda$;
                  \item the algebraic dimension of $P_{\delta+1}/P_\delta$ is
                        at most countable for every $\delta<\lambda$.
              \end{enumerate}
    \end{enumerate}
\end{theorem}
\begin{proof} $(i)\iff(ii)$ follows from theorem~\ref{FinalCharacterization},
    since the completion of $P$ is $\ellone{\Gamma}$.

    $(iii)\implies(i)$ Let $0<\theta<1/2$. By the remark after
    theorem~\ref{TransfiniteBackAndForth}, conditions $(a)$ and $(c)$ imply
    condition $(iii)$ of that theorem for the given filtration and $\theta$,
    while $(a)$ and $(b)$ are its conditions $(i)$ and $(ii)$. Hence there
    is a linear bijection $T$ of $\ellone{\Gamma}$ with $T(P)=\lo{\Gamma}$
    and $\norm{T}\norm{T^{-1}}\leq{(1+\theta)}^2{(1-\theta)}^{-1}$. As
    $\theta$ is arbitrary, $\BM(P,\lo{\Gamma})=1$, and $P$ is extremely
    projective by theorem~\ref{FinalCharacterization}.

    $(ii)\implies(iii)$ Fix $S$, $\sigma$, $\pi$ and the sequences
    $p^\gamma_n$ as above, and an enumeration ${(\gamma_\alpha)}_{\alpha<\kappa}$
    of $\Gamma$ by an ordinal $\kappa$. By transfinite recursion we define an
    increasing family ${((A_\alpha,B_\alpha))}_{\alpha\leq\kappa}$ of closed
    pairs: $A_0=B_0=\varnothing$, which is a closed pair since
    $\sigma(0)=0$; given $(A_\alpha,B_\alpha)$ with $\alpha<\kappa$, let
    $(A_{\alpha+1},B_{\alpha+1})$ be a closed pair provided by
    lemma~\ref{CountableClosingOff} for $(A_\alpha,B_\alpha)$ and
    $N=\{\gamma_\alpha\}$, so that $A_{\alpha+1}\setminus A_\alpha$ is at
    most countable; for a limit ordinal $\alpha\leq\kappa$ let $A_\alpha$ and
    $B_\alpha$ be the unions of the previous sets, which form a closed pair
    by lemma~\ref{UnionsOfClosedPairs}. Since $\gamma_\alpha\in A_{\alpha+1}$,
    we have $A_\kappa=\Gamma$, so ${(A_\alpha)}_{\alpha\leq\kappa}$ is a
    filtration of $\Gamma$ with at most countable blocks. Condition $(a)$
    holds by lemma~\ref{ClosedPairProps}, condition $(b)$ by
    lemma~\ref{UnionsOfClosedPairs}, and condition $(c)$ by
    lemma~\ref{CountableQuotient}.
\end{proof}

Condition $(iii)$ of theorem~\ref{CountableBlocksCriterion} refers only to the
subspaces $P\cap\ellone{A}$, $A\subset\Gamma$. For sets $\Gamma$ of
cardinality $\aleph_1$ it can be stated as a condition on a single chain of
subspaces. In the rest of this section we take $\Gamma=\omega_1$; every set
of cardinality $\aleph_1$ can be enumerated by $\omega_1$, and we shall see in
proposition~\ref{InvarianceOfDelta} that the choice of the enumeration is
immaterial. For $\delta\leq\omega_1$ we write $P_\delta=P\cap\ellone{\delta}$,
where $\ellone{\delta}$ is the subspace of $\ellone{\omega_1}$ consisting of
the families supported in $\delta=\{\beta:\beta<\delta\}$. We shall use the
following standard facts about closed unbounded subsets of $\omega_1$
\cite[chapter 8]{JechSetTheory}: the intersection of two closed unbounded
sets is closed unbounded, and for every function $h:\omega_1\to\omega_1$ the
set of ordinals $\delta<\omega_1$ such that $h(\alpha)<\delta$ for all
$\alpha<\delta$ is closed unbounded; we call such ordinals closed under $h$.
Moreover, the increasing enumeration ${(c_\alpha)}_{\alpha<\omega_1}$ of a
closed unbounded set $C$ is continuous, that is,
$c_\alpha=\sup_{\alpha'<\alpha}c_{\alpha'}$ for every limit ordinal
$\alpha<\omega_1$.

\begin{definition}\label{DefDelta} Let $P$ be a dense subspace of
    $\ellone{\omega_1}$. We denote by $\Delta(P)$ the set of limit ordinals
    $\delta<\omega_1$ such that $P_\delta\neq\bigcup_{\delta'<\delta}P_{\delta'}$.
\end{definition}

\begin{theorem}\label{AlephOneCriterion} Let $P$ be a dense subspace of
    $\ellone{\omega_1}$. Then $P$ is extremely projective if and only if
    \begin{enumerate}[label = (\roman*)]
        \item $P_\delta$ has at most countable algebraic dimension for every
              $\delta<\omega_1$, and
        \item the set $\Delta(P)$ is not stationary in $\omega_1$.
    \end{enumerate}
    Condition $(i)$ is equivalent to the requirement that every separable
    subspace of $P$ has at most countable algebraic dimension.
\end{theorem}
\begin{proof} First we prove the last claim. Each $P_\delta$, $\delta<\omega_1$,
    is separable as a subspace of $\ellone{\delta}$. Conversely, if $Z$ is a
    separable subspace of $P$, then the supports of the elements of a
    countable dense subset of $Z$ are contained in a countable subset of
    $\omega_1$, hence in some $\delta<\omega_1$, and by the continuity of the
    coordinate functionals $Z\subset P_\delta$.

    Let $P$ be extremely projective. Condition $(i)$ holds by
    theorem~\ref{SeparableSubspaces}. To verify $(ii)$ we fix $S$, $\sigma$,
    $\pi$ and $p^\gamma_n$ as in section~\ref{SectionUncountable} and apply
    the construction of the proof of theorem~\ref{CountableBlocksCriterion}
    to the enumeration $\gamma_\alpha=\alpha$ of $\omega_1$. We obtain an
    increasing family ${((A_\alpha,B_\alpha))}_{\alpha\leq\omega_1}$ of closed
    pairs, continuous at limit ordinals, such that $A_\alpha$ is countable
    for $\alpha<\omega_1$, $\alpha\in A_{\alpha+1}$ and $A_{\omega_1}=\omega_1$.
    Define $f,g:\omega_1\to\omega_1$ by $f(\alpha)=\sup A_\alpha+1$ and
    $g(\beta)=\min\{\alpha:\beta\in A_\alpha\}$, and let $D$ be the set of
    limit ordinals $\delta<\omega_1$ closed under $f$ and $g$; it is closed
    unbounded. Let $\delta\in D$. Then
    $A_\delta=\bigcup_{\alpha<\delta}A_\alpha\subset\delta$, because
    $\sup A_\alpha<f(\alpha)<\delta$ for $\alpha<\delta$, and
    $\delta\subset A_\delta$, because $\beta<\delta$ implies
    $\beta\in A_{g(\beta)}\subset A_\delta$. Thus, $A_\delta=\delta$, and by
    lemma~\ref{UnionsOfClosedPairs}
    \[
    P_\delta=P_{A_\delta}=\bigcup_{\alpha<\delta}P_{A_\alpha}
    \subset\bigcup_{\alpha<\delta}P_{f(\alpha)}
    \subset\bigcup_{\delta'<\delta}P_{\delta'}\subset P_\delta,
    \]
    since $A_\alpha\subset f(\alpha)<\delta$. Hence $\delta\notin\Delta(P)$.
    So $\Delta(P)\cap D=\varnothing$, and $\Delta(P)$ is not stationary.

    Conversely, assume $(i)$ and $(ii)$. For $\gamma<\omega_1$ fix a sequence
    ${(p^\gamma_n)}_{n\in\mathbb{N}}$ in $P$ converging to $e_\gamma$ and put
    $h(\gamma)=\sup\bigcup_n\supp(p^\gamma_n)+1<\omega_1$. Let $C_1$ be the
    set of ordinals $\delta<\omega_1$ closed under $h$; for $\delta\in C_1$
    the space $P_\delta$ contains $p^\gamma_n$ for all $\gamma<\delta$, so
    $P_\delta$ is dense in $\ellone{\delta}$. By $(ii)$ there is a closed
    unbounded set $C_2$ disjoint from $\Delta(P)$. Let $C=(C_1\cap C_2)\cup\{0\}$,
    a closed unbounded set, and let ${(c_\alpha)}_{\alpha<\omega_1}$ be its
    increasing enumeration, so that $c_0=0$. Put $\Gamma_\alpha=c_\alpha$ for
    $\alpha<\omega_1$ and $\Gamma_{\omega_1}=\omega_1$. Then
    ${(\Gamma_\alpha)}_{\alpha\leq\omega_1}$ is a filtration of $\omega_1$:
    it is increasing, $\Gamma_0=\varnothing$, continuity at limit ordinals
    $\alpha<\omega_1$ is the continuity of the enumeration of $C$, and
    continuity at $\omega_1$ holds because $C$ is unbounded. Its blocks
    $N_\alpha=\{\beta:c_\alpha\leq\beta<c_{\alpha+1}\}$ are countable. We
    verify conditions $(a)$--$(c)$ of theorem~\ref{CountableBlocksCriterion}.
    Condition $(a)$ holds for $\alpha<\omega_1$ because $c_\alpha\in C_1$ (for
    $\alpha=0$ trivially), and for $\alpha=\omega_1$ because $P$ is dense.
    For a limit ordinal $\alpha<\omega_1$ the ordinal $c_\alpha$ is a limit
    ordinal in $C_2$, hence $c_\alpha\notin\Delta(P)$, and since the ordinals
    $c_{\alpha'}$, $\alpha'<\alpha$, are cofinal in $c_\alpha$ and the
    spaces $P_\delta$ increase with $\delta$,
    \[
    P_{c_\alpha}=\bigcup_{\delta'<c_\alpha}P_{\delta'}
    =\bigcup_{\alpha'<\alpha}P_{c_{\alpha'}};
    \]
    at $\alpha=\omega_1$ condition $(b)$ holds because every element of $P$
    has countable support and therefore belongs to $P_{c_{\alpha'}}$ for some
    $\alpha'<\omega_1$. Finally, $(c)$ follows from $(i)$, since
    $P_{c_{\alpha+1}}$ has at most countable algebraic dimension. By
    theorem~\ref{CountableBlocksCriterion} the space $P$ is extremely
    projective.
\end{proof}

Neither of the conditions of theorem~\ref{AlephOneCriterion} can be dropped:
the ladder space $P_L$ of example~\ref{LadderSpace} satisfies $(i)$ but not
$(ii)$, as we shall see again below, while the space
$P=P_w+\lo{\omega_1\setminus\omega}$, where $P_w\subset\ellone{\omega}$ is the
weighted space of example~\ref{WeightedExample}, satisfies $(ii)$ but not
$(i)$: here $\Delta(P)=\{\omega\}$, because every element of $P$ is the sum
of an element of $P_w$ and a finitely supported family, so the chain
${(P_\delta)}_\delta$ is continuous at every limit ordinal $\delta>\omega$,
while $P_\omega=P_w$ has uncountable algebraic dimension.

The set $\Delta(P)$ depends on the position of $P$ in $\ellone{\omega_1}$,
but only up to a non-stationary set.

\begin{proposition}\label{InvarianceOfDelta} Let $P$ and $P'$ be dense
    subspaces of $\ellone{\omega_1}$ and $T:\ellone{\omega_1}\to\ellone{\omega_1}$
    be a bounded linear bijection with $T(P)=P'$. Then
    $\Delta(P)\triangle\Delta(P')$ is not stationary. In particular, if $P$
    and $P'$ are topologically isomorphic as normed spaces, then
    $\Delta(P)\triangle\Delta(P')$ is not stationary, and the class of
    $\Delta(P)$ modulo non-stationary sets does not depend on the enumeration
    of the index set by $\omega_1$.
\end{proposition}
\begin{proof} For $\delta<\omega_1$ put
    \[
    F(\delta)=\sup_{\beta<\delta}\bigl(\sup\supp(T(e_\beta))+1\bigr),\qquad
    G(\delta)=\sup_{\beta<\delta}\bigl(\sup\supp(T^{-1}(e_\beta))+1\bigr),
    \]
    which are countable ordinals, since all supports are countable. Let $D$
    be the set of limit ordinals $\delta<\omega_1$ closed under $F$ and $G$;
    it is closed unbounded. Let $\delta\in D$ and $\delta'<\delta$. Then
    $T(e_\beta)\in\ellone{F(\delta')}$ for all $\beta<\delta'$, so by the
    continuity of $T$ we have $T(\ellone{\delta'})\subset\ellone{F(\delta')}
    \subset\ellone{\delta}$, and since $\ellone{\delta}$ is the closure of
    $\bigcup_{\delta'<\delta}\ellone{\delta'}$, also
    $T(\ellone{\delta})\subset\ellone{\delta}$. The same holds for $T^{-1}$
    with $G$ in place of $F$, so $T(\ellone{\delta})=\ellone{\delta}$ and
    $T(P_\delta)=T(P)\cap T(\ellone{\delta})=P'_\delta$. Moreover,
    $T(P_{\delta'})\subset P'\cap\ellone{F(\delta')}=P'_{F(\delta')}$ with
    $F(\delta')<\delta$, so
    $T\bigl(\bigcup_{\delta'<\delta}P_{\delta'}\bigr)\subset
    \bigcup_{\delta'<\delta}P'_{\delta'}$, and symmetrically for $T^{-1}$;
    hence $T$ maps $\bigcup_{\delta'<\delta}P_{\delta'}$ onto
    $\bigcup_{\delta'<\delta}P'_{\delta'}$. Consequently, $\delta\in\Delta(P)$
    if and only if $\delta\in\Delta(P')$, for every $\delta\in D$. Thus,
    $\Delta(P)\triangle\Delta(P')$ is disjoint from $D$ and is not
    stationary.

    If $P$ and $P'$ are topologically isomorphic, then an isomorphism
    $P\to P'$ extends to a bounded linear bijection of the completions, which
    are both $\ellone{\omega_1}$, so the first claim applies. Finally, two
    enumerations of the same index set by $\omega_1$ differ by a permutation
    of $\omega_1$, which induces an isometric bijection of
    $\ellone{\omega_1}$ mapping $P$ onto the same space $P$ regarded with
    respect to the other enumeration.
\end{proof}

We now describe the ladder spaces built on arbitrary sets of limit ordinals.
Let $E$ be a set of limit ordinals below $\omega_1$. For every $\alpha\in E$
fix a strictly increasing sequence ${(\alpha_n)}_{n\in\mathbb{N}}$ with
$\sup_n\alpha_n=\alpha$ and put $v_\alpha=\sum_{n=1}^\infty 2^{-n}e_{\alpha_n}$,
as in example~\ref{LadderSpace}, and let
\[
P_E=\lo{\omega_1}+\Lin\{v_\alpha:\alpha\in E\}\subset\ellone{\omega_1}.
\]
Thus, $P_L=P_E$ for the set $E$ of all limit ordinals below $\omega_1$.

\begin{theorem}\label{LadderSpacesOnE} Let $E$ and $E'$ be sets of limit
    ordinals below $\omega_1$. Then
    \begin{enumerate}[label = (\roman*)]
        \item $P_E$ is a dense subspace of $\ellone{\omega_1}$ of algebraic
              dimension $\aleph_1$, and for every $\delta<\omega_1$ the space
              $(P_E)_\delta=\lo{\delta}+\Lin\{v_\alpha:\alpha\in E,\ \alpha\leq\delta\}$
              has countable algebraic dimension and is extremely projective;
        \item $\Delta(P_E)=E$;
        \item $P_E$ is extremely projective if and only if $E$ is not
              stationary;
        \item if $E\triangle E'$ is stationary, then $P_E$ and $P_{E'}$ are
              not topologically isomorphic.
    \end{enumerate}
    Consequently, there are $2^{\aleph_1}$ pairwise non-isomorphic dense
    subspaces of $\ellone{\omega_1}$ of algebraic dimension $\aleph_1$ all of
    whose subspaces $P\cap\ellone{\delta}$, $\delta<\omega_1$, are extremely
    projective, and only one of them, up to a topological isomorphism, is
    extremely projective itself.
\end{theorem}
\begin{proof} $(i)$ The space $P_E$ contains $\lo{\omega_1}$, so it is dense
    and has algebraic dimension $\aleph_1$. The proof of part $(ii)$ of
    example~\ref{LadderSpace} uses only the fact that
    $\supp(v_\alpha)\cap\supp(v_\beta)$ is finite for $\alpha<\beta$, which
    holds for all limit ordinals $\alpha<\beta$ regardless of $E$; it gives
    the formula for $(P_E)_\delta$, and this space has countable algebraic
    dimension and is dense in $\ellone{\delta}$, hence extremely projective
    by theorem~\ref{MainTheorem}.

    $(ii)$ Let $\delta<\omega_1$ be a limit ordinal. By $(i)$,
    \[
    \bigcup_{\delta'<\delta}(P_E)_{\delta'}=\lo{\delta}+\Lin\{v_\alpha:\alpha\in E,\ \alpha<\delta\},
    \]
    because $\lo{\delta}=\bigcup_{\delta'<\delta}\lo{\delta'}$ and
    $v_\alpha\in(P_E)_\alpha$ for $\alpha\in E$. If $\delta\notin E$, this is
    $(P_E)_\delta$ by $(i)$, so $\delta\notin\Delta(P_E)$. If $\delta\in E$,
    then $v_\delta\in(P_E)_\delta$, while
    $v_\delta\notin\lo{\delta}+\Lin\{v_\alpha:\alpha\in E,\ \alpha<\delta\}$:
    if $v_\delta=y+\sum_{\alpha\in F}c_\alpha v_\alpha$ with $y\in\lo{\delta}$
    and a finite set $F$ of limit ordinals below $\delta$, then for all
    sufficiently large $n$ the coordinate of the right hand side at
    $\delta_n$ vanishes, because $\supp(y)$ and $\supp(v_\alpha)\cap\supp(v_\delta)$,
    $\alpha\in F$, are finite, whereas the coordinate of $v_\delta$ at
    $\delta_n$ is $2^{-n}$. Hence $\delta\in\Delta(P_E)$.

    $(iii)$ follows from $(i)$, $(ii)$ and theorem~\ref{AlephOneCriterion}, and
    $(iv)$ follows from $(ii)$ and proposition~\ref{InvarianceOfDelta}.

    For the last claim recall that $\omega_1$ is the union of $\aleph_1$
    pairwise disjoint stationary sets ${(S_\xi)}_{\xi<\omega_1}$
    \cite[theorem 8.10]{JechSetTheory}; replacing $S_\xi$ by its intersection
    with the closed unbounded set of limit ordinals, we may assume that the
    sets $S_\xi$ consist of limit ordinals. For $X\subset\omega_1$ put
    $E_X=\bigcup_{\xi\in X}S_\xi$. If $X\neq Y$, then $E_X\triangle E_Y$
    contains some $S_\xi$ and is stationary, so by $(iv)$ the spaces
    $P_{E_X}$, $X\subset\omega_1$, are pairwise non-isomorphic. By $(iii)$,
    $P_{E_X}$ is extremely projective only for $X=\varnothing$, and every
    extremely projective normed space of algebraic dimension $\aleph_1$ is
    topologically isomorphic to $\lo{\omega_1}$ by
    theorem~\ref{FinalCharacterization}.
\end{proof}

\begin{remark}\label{RemarkAlmostFree} Theorems~\ref{AlephOneCriterion}
    and~\ref{LadderSpacesOnE} give a second proof of part $(iii)$ of
    example~\ref{LadderSpace} which does not use Fodor's lemma; the latter
    is now needed only for theorem~\ref{FodorObstruction} and
    corollary~\ref{OmegaOneContinuity}, which concern arbitrary
    $\omega_1$-sequences of vectors rather than the coordinate subspaces
    $P_\delta$. The results of this section are parallel to the theory of
    almost free abelian groups, see \cite{EklofMeklerAlmostFree}:
    theorem~\ref{CountableBlocksCriterion} corresponds to the criteria of
    Pontryagin and Hill for freeness in terms of filtrations by countable
    subgroups with free quotients, the class of $\Delta(P)$ modulo
    non-stationary sets corresponds to the invariant $\Gamma(G)$ of Eklof,
    and the spaces $P_E$ correspond to the $\aleph_1$-free groups of
    cardinality $\aleph_1$ built from ladder systems. In that theory the
    situation for cardinals above $\aleph_1$ depends on set-theoretic
    hypotheses, and we expect the same here. We do not know whether
    condition $(iii)$ of theorem~\ref{CountableBlocksCriterion} can be
    replaced, for $|\Gamma|\geq\aleph_2$, by a condition on the subspaces
    $P\cap\ellone{A}$ for $A$ ranging over a closed unbounded family of
    subsets of $\Gamma$ of cardinality less than $|\Gamma|$, nor whether
    there is a dense subspace of $\ellone{\omega_2}$ which is not
    topologically projective although $P\cap\ellone{A}$ is extremely
    projective for every $A\subset\omega_2$ of cardinality at most $\aleph_1$.
\end{remark}

\begin{thebibliography}{999}
    %
    \bibitem{HelHomolBanTopAlg}
    \textit{A. Ya. Helemskii} The homology of Banach and topological algebras,
    Kluwer, Dordrecht (1989)
    %
    %
    \bibitem{HelMetrFrQMod}
    \textit{A. Ya. Helemskii} Metric freeness and projectivity for classical
    and quantum normed modules, Sb. Math., 204:7 (2013), 1056--1083
    %
    %
    \bibitem{HelExtrProjSeqAlg}
    \textit{A. Ya. Helemskii} Extreme version of projectivity for normed
    modules over sequence algebras, Canad. J. Math., 65:3 (2013), 559--574;
    arXiv:1104.2463 (the numbering of statements refers to the arXiv
    version, titled ``Metric version of projectivity for normed modules over
    sequence algebras'')
    %
    %
    \bibitem{WhiteInjmoduAlg}
    \textit{M. C. White} Injective modules for uniform algebras, Proc. London
    Math. Soc. (3), 73:1 (1996), 155--184
    %
    %
    \bibitem{KotheTopProjBanSp}
    \textit{G. K\"othe} Hebbare lokalkonvexe R\"aume, Math. Ann., 165 (1966),
    181--195
    %
    %
    \bibitem{GrothMetrProjFlatBanSp}
    \textit{A. Grothendieck} Une caract\'erisation vectorielle-m\'etrique des
    espaces $L_1$, Canad. J. Math., 7 (1955), 552--561
    %
    %
    \bibitem{LaceyIsomThOfClassicBanSp}
    \textit{H. E. Lacey} The isometric theory of classical Banach spaces,
    Springer, Berlin (1974)
    %
    %
    \bibitem{GronLiftProbNormSp}
    \textit{N. Gr{\o}nb{\ae}k} Lifting problems for normed spaces, Ann. Funct.
    Anal., 7:1 (2016), 118--126
    %
    %
    \bibitem{HajekRussoDenselyIsom}
    \textit{P. H\'ajek, T. Russo} On densely isomorphic normed spaces,
    J. Funct. Anal., 279:11 (2020), 108773; arXiv:1910.01527
    %
    %
    \bibitem{AvilesKoszmiderLadder}
    \textit{A. Avil\'es, P. Koszmider} Almost disjoint families of countable
    sets and separable complementation properties, J. Math. Anal. Appl.,
    405:1 (2013), 252--263; arXiv:1209.0199
    %
    %
    \bibitem{LowigDimLinRaume}
    \textit{H. L\"owig} \"Uber die Dimension linearer R\"aume, Studia Math.,
    5 (1934), 18--23
    %
    %
    \bibitem{JechSetTheory}
    \textit{T. Jech} Set theory. The third millennium edition, Springer,
    Berlin (2003)
    %
    %
    \bibitem{EklofMeklerAlmostFree}
    \textit{P. C. Eklof, A. H. Mekler} Almost free modules. Set-theoretic
    methods, revised edition, North-Holland Mathematical Library, 65,
    North-Holland, Amsterdam (2002)
    %
\end{thebibliography}

Norbert Nemesh, Faculty of Mechanics and Mathematics, Moscow State University,
Moscow 119991 Russia

\textit{E-mail address:} nemeshnorbert@yandex.ru

\end{document}
